%%%%%%%%%%%%%%%%%%%%%%%%%%%%%%%%%%%%%%%%%%%%%%%%%%%%%%%%%%%%
% 11.7 SCIENTIFIC COMPUTING
% The Composition of Advanced Mathematics
%%%%%%%%%%%%%%%%%%%%%%%%%%%%%%%%%%%%%%%%%%%%%%%%%%%%%%%%%%%%

\section{Scientific Computing}
\label{sec:scientific-computing}

\prerequisite{
Numerical analysis, numerical linear algebra, numerical optimization,
numerical differential equations, Monte Carlo methods, basic algorithms,
programming concepts, and elementary computer architecture.
}

\learningobjectives{
By the end of this section, the reader should be able to:
\begin{itemize}
    \item explain the role of scientific computing in applied mathematics;
    \item distinguish mathematical models, numerical algorithms, and
          software implementations;
    \item understand the scientific-computing workflow;
    \item recognize sources of computational error;
    \item distinguish algorithmic complexity from practical runtime;
    \item understand arithmetic intensity conceptually;
    \item recognize memory and data-movement costs;
    \item understand cache locality conceptually;
    \item distinguish scalar and vectorized computation;
    \item understand matrix-oriented computation;
    \item recognize sparse data structures;
    \item understand the importance of numerical libraries;
    \item distinguish interpreted and compiled execution conceptually;
    \item recognize profiling and benchmarking;
    \item understand reproducible computational experiments;
    \item understand deterministic and nondeterministic computation;
    \item recognize parallel computing;
    \item distinguish shared-memory and distributed-memory parallelism;
    \item understand task and data parallelism;
    \item recognize multicore computing;
    \item recognize vector processors and GPUs;
    \item understand distributed-memory communication conceptually;
    \item recognize message passing;
    \item understand speedup and parallel efficiency;
    \item apply Amdahl's law;
    \item understand weak and strong scaling;
    \item recognize load balancing;
    \item understand communication overhead;
    \item recognize synchronization costs;
    \item understand domain decomposition;
    \item recognize iterative methods suited to parallel computation;
    \item understand precision choices;
    \item distinguish single, double, and mixed precision;
    \item recognize arbitrary-precision arithmetic;
    \item understand reproducibility limitations in floating-point parallel
          reductions;
    \item understand software verification and numerical validation;
    \item recognize unit testing for mathematical software;
    \item understand convergence testing;
    \item understand dimensional and unit checks;
    \item recognize data provenance;
    \item understand checkpointing;
    \item recognize restartable computations;
    \item understand parameter sweeps;
    \item recognize ensemble simulation;
    \item understand computational pipelines;
    \item recognize workflow automation;
    \item understand scientific visualization conceptually;
    \item distinguish visualization from quantitative validation;
    \item recognize high-performance computing;
    \item understand computational scalability;
    \item recognize accelerator computing;
    \item understand performance portability conceptually;
    \item connect scientific computing with simulation, optimization,
          statistics, data analysis, differential equations, and applied
          mathematical modeling.
\end{itemize}
}

%%%%%%%%%%%%%%%%%%%%%%%%%%%%%%%%%%%%%%%%%%%%%%%%%%%%%%%%%%%%
\subsection{What Is Scientific Computing?}
\label{subsec:what-is-scientific-computing}
%%%%%%%%%%%%%%%%%%%%%%%%%%%%%%%%%%%%%%%%%%%%%%%%%%%%%%%%%%%%

Scientific computing uses mathematical models, numerical algorithms, and
computer systems to study problems through computation.

A typical scientific-computing pipeline is

\[
\boxed{
\text{mathematical model}
\longrightarrow
\text{numerical method}
\longrightarrow
\text{implementation}
\longrightarrow
\text{computed result}.
}
\]

Each stage can introduce its own assumptions and errors.

%%%%%%%%%%%%%%%%%%%%%%%%%%%%%%%%%%%%%%%%%%%%%%%%%%%%%%%%%%%%
\subsection{Scientific Computing as Applied Mathematics}
\label{subsec:scientific-computing-applied-mathematics}
%%%%%%%%%%%%%%%%%%%%%%%%%%%%%%%%%%%%%%%%%%%%%%%%%%%%%%%%%%%%

Scientific computing is not merely programming.

It combines:

\[
\boxed{
\text{mathematical modeling},
}
\]

\[
\boxed{
\text{numerical analysis},
}
\]

\[
\boxed{
\text{algorithm design},
}
\]

and

\[
\boxed{
\text{computer implementation}.
}
\]

The mathematical quality of the result depends on all four.

%%%%%%%%%%%%%%%%%%%%%%%%%%%%%%%%%%%%%%%%%%%%%%%%%%%%%%%%%%%%
\subsection{Model, Method, and Implementation}
\label{subsec:model-method-implementation}
%%%%%%%%%%%%%%%%%%%%%%%%%%%%%%%%%%%%%%%%%%%%%%%%%%%%%%%%%%%%

Three layers should be distinguished carefully.

The \textbf{model} specifies the mathematical problem:

\[
\boxed{
F(u,\theta)=0.
}
\]

The \textbf{numerical method} approximates the solution:

\[
\boxed{
F_h(u_h,\theta)=0.
}
\]

The \textbf{implementation} realizes the numerical method as software.

A result can fail because of an error at any of these layers.

%%%%%%%%%%%%%%%%%%%%%%%%%%%%%%%%%%%%%%%%%%%%%%%%%%%%%%%%%%%%
\subsection{Sources of Computational Error}
\label{subsec:scientific-computing-error-sources}
%%%%%%%%%%%%%%%%%%%%%%%%%%%%%%%%%%%%%%%%%%%%%%%%%%%%%%%%%%%%

Important error sources include:

\[
\boxed{
\text{modeling error},
}
\]

\[
\boxed{
\text{data uncertainty},
}
\]

\[
\boxed{
\text{discretization error},
}
\]

\[
\boxed{
\text{iteration error},
}
\]

\[
\boxed{
\text{floating-point error},
}
\]

and

\[
\boxed{
\text{software error}.
}
\]

These should be analyzed separately whenever possible.

%%%%%%%%%%%%%%%%%%%%%%%%%%%%%%%%%%%%%%%%%%%%%%%%%%%%%%%%%%%%
\subsection{Computational Workflow}
\label{subsec:scientific-computing-workflow}
%%%%%%%%%%%%%%%%%%%%%%%%%%%%%%%%%%%%%%%%%%%%%%%%%%%%%%%%%%%%

A practical workflow is:

\begin{enumerate}
    \item define the scientific question;
    \item construct the mathematical model;
    \item identify inputs, parameters, and outputs;
    \item determine the numerical problem class;
    \item choose appropriate numerical methods;
    \item estimate expected computational cost;
    \item implement or select trusted software;
    \item verify individual components;
    \item run benchmark problems;
    \item perform convergence studies;
    \item perform sensitivity studies;
    \item evaluate computational performance;
    \item document configuration and software versions;
    \item store sufficient information for reproducibility;
    \item interpret numerical results within model limitations.
\end{enumerate}

%%%%%%%%%%%%%%%%%%%%%%%%%%%%%%%%%%%%%%%%%%%%%%%%%%%%%%%%%%%%
\subsection{Algorithmic Complexity}
\label{subsec:scientific-computing-complexity}
%%%%%%%%%%%%%%%%%%%%%%%%%%%%%%%%%%%%%%%%%%%%%%%%%%%%%%%%%%%%

Algorithmic complexity estimates computational work as problem size grows.

Examples include:

\[
\boxed{
O(n),
}
\]

\[
\boxed{
O(n\log n),
}
\]

\[
\boxed{
O(n^2),
}
\]

and

\[
\boxed{
O(n^3).
}
\]

Complexity describes asymptotic growth rather than exact execution time.

%%%%%%%%%%%%%%%%%%%%%%%%%%%%%%%%%%%%%%%%%%%%%%%%%%%%%%%%%%%%
\subsection{Operation Count}
\label{subsec:operation-count}
%%%%%%%%%%%%%%%%%%%%%%%%%%%%%%%%%%%%%%%%%%%%%%%%%%%%%%%%%%%%

Numerical algorithms are often analyzed by counting arithmetic operations.

For dense Gaussian elimination, the leading arithmetic work is

\[
\boxed{
\frac{2}{3}n^3.
}
\]

For a dense matrix-vector product,

\[
\boxed{
A x,
}
\]

the work is approximately

\[
\boxed{
O(n^2).
}
\]

Operation counts are useful but do not completely predict modern hardware
performance.

%%%%%%%%%%%%%%%%%%%%%%%%%%%%%%%%%%%%%%%%%%%%%%%%%%%%%%%%%%%%
\subsection{Runtime Is More Than Arithmetic}
\label{subsec:runtime-more-than-arithmetic}
%%%%%%%%%%%%%%%%%%%%%%%%%%%%%%%%%%%%%%%%%%%%%%%%%%%%%%%%%%%%

Actual runtime depends on:

\[
\boxed{
\text{arithmetic operations},
}
\]

\[
\boxed{
\text{memory access},
}
\]

\[
\boxed{
\text{communication},
}
\]

\[
\boxed{
\text{synchronization},
}
\]

and

\[
\boxed{
\text{hardware utilization}.
}
\]

A method with fewer arithmetic operations can still be slower if it moves
data inefficiently.

%%%%%%%%%%%%%%%%%%%%%%%%%%%%%%%%%%%%%%%%%%%%%%%%%%%%%%%%%%%%
\subsection{Memory Hierarchy}
\label{subsec:memory-hierarchy}
%%%%%%%%%%%%%%%%%%%%%%%%%%%%%%%%%%%%%%%%%%%%%%%%%%%%%%%%%%%%

Modern computers contain multiple levels of memory.

Conceptually,

\[
\boxed{
\text{registers}
\rightarrow
\text{cache}
\rightarrow
\text{main memory}
\rightarrow
\text{storage}.
}
\]

Memory closer to the processor is typically faster but smaller.

Efficient scientific software attempts to reuse data before it is moved
again from slower memory.

%%%%%%%%%%%%%%%%%%%%%%%%%%%%%%%%%%%%%%%%%%%%%%%%%%%%%%%%%%%%
\subsection{Data Locality}
\label{subsec:data-locality}
%%%%%%%%%%%%%%%%%%%%%%%%%%%%%%%%%%%%%%%%%%%%%%%%%%%%%%%%%%%%

Temporal locality means recently accessed data are likely to be used
again.

Spatial locality means nearby memory locations are likely to be used
together.

Algorithms that exploit locality can significantly outperform
mathematically equivalent implementations.

%%%%%%%%%%%%%%%%%%%%%%%%%%%%%%%%%%%%%%%%%%%%%%%%%%%%%%%%%%%%
\subsection{Arithmetic Intensity}
\label{subsec:arithmetic-intensity}
%%%%%%%%%%%%%%%%%%%%%%%%%%%%%%%%%%%%%%%%%%%%%%%%%%%%%%%%%%%%

Arithmetic intensity is approximately

\[
\boxed{
\frac{
\text{floating-point operations}
}{
\text{bytes moved}
}.
}
\]

High arithmetic intensity means substantial computation is performed for
each unit of data movement.

Low arithmetic intensity indicates that memory bandwidth may dominate
performance.

%%%%%%%%%%%%%%%%%%%%%%%%%%%%%%%%%%%%%%%%%%%%%%%%%%%%%%%%%%%%
\subsection{Memory-Bound and Compute-Bound Algorithms}
\label{subsec:memory-compute-bound}
%%%%%%%%%%%%%%%%%%%%%%%%%%%%%%%%%%%%%%%%%%%%%%%%%%%%%%%%%%%%

A memory-bound algorithm spends much of its time moving data.

A compute-bound algorithm spends much of its time performing arithmetic.

For example, sparse matrix-vector multiplication is often memory-bound
because relatively little arithmetic is performed per stored matrix
entry.

Dense matrix multiplication has much higher arithmetic intensity.

%%%%%%%%%%%%%%%%%%%%%%%%%%%%%%%%%%%%%%%%%%%%%%%%%%%%%%%%%%%%
\subsection{Vectorization}
\label{subsec:scientific-vectorization}
%%%%%%%%%%%%%%%%%%%%%%%%%%%%%%%%%%%%%%%%%%%%%%%%%%%%%%%%%%%%

Vectorized computation applies operations to collections of values rather
than processing one scalar element at a time at the high-level programming
layer.

For example,

\[
\boxed{
y
=
Ax+b
}
\]

expresses a complete matrix-vector operation.

Well-designed numerical libraries can execute such operations efficiently
using optimized low-level routines.

%%%%%%%%%%%%%%%%%%%%%%%%%%%%%%%%%%%%%%%%%%%%%%%%%%%%%%%%%%%%
\subsection{Matrix-Oriented Computation}
\label{subsec:matrix-oriented-computation}
%%%%%%%%%%%%%%%%%%%%%%%%%%%%%%%%%%%%%%%%%%%%%%%%%%%%%%%%%%%%

Many scientific problems can be expressed using:

\[
\boxed{
\text{vectors},
}
\]

\[
\boxed{
\text{matrices},
}
\]

and

\[
\boxed{
\text{tensor operations}.
}
\]

This provides both mathematical clarity and access to highly optimized
numerical kernels.

%%%%%%%%%%%%%%%%%%%%%%%%%%%%%%%%%%%%%%%%%%%%%%%%%%%%%%%%%%%%
\subsection{Numerical Libraries}
\label{subsec:numerical-libraries}
%%%%%%%%%%%%%%%%%%%%%%%%%%%%%%%%%%%%%%%%%%%%%%%%%%%%%%%%%%%%

Scientific software frequently relies on established numerical libraries
for operations such as:

\[
\boxed{
\text{matrix multiplication},
}
\]

\[
\boxed{
\text{factorization},
}
\]

\[
\boxed{
\text{eigenvalue computation},
}
\]

and

\[
\boxed{
\text{Fourier transforms}.
}
\]

Using a tested optimized library is often preferable to reimplementing
fundamental numerical kernels.

%%%%%%%%%%%%%%%%%%%%%%%%%%%%%%%%%%%%%%%%%%%%%%%%%%%%%%%%%%%%
\subsection{Layered Numerical Software}
\label{subsec:layered-numerical-software}
%%%%%%%%%%%%%%%%%%%%%%%%%%%%%%%%%%%%%%%%%%%%%%%%%%%%%%%%%%%%

A high-level solver may call lower-level routines.

For example,

\[
\boxed{
\text{optimization algorithm}
}
\]

may call

\[
\boxed{
\text{sparse linear solver},
}
\]

which calls

\[
\boxed{
\text{matrix-vector kernels}.
}
\]

Understanding these layers helps locate both performance bottlenecks and
numerical failures.

%%%%%%%%%%%%%%%%%%%%%%%%%%%%%%%%%%%%%%%%%%%%%%%%%%%%%%%%%%%%
\subsection{Compiled and Interpreted Execution}
\label{subsec:compiled-interpreted}
%%%%%%%%%%%%%%%%%%%%%%%%%%%%%%%%%%%%%%%%%%%%%%%%%%%%%%%%%%%%

Compiled languages translate programs into machine-level instructions
before execution.

Interpreted or dynamically executed environments may evaluate higher-level
operations during runtime.

For scientific computing, the distinction is often less important than
whether expensive numerical operations are delegated to optimized compiled
libraries.

%%%%%%%%%%%%%%%%%%%%%%%%%%%%%%%%%%%%%%%%%%%%%%%%%%%%%%%%%%%%
\subsection{Profiling}
\label{subsec:scientific-profiling}
%%%%%%%%%%%%%%%%%%%%%%%%%%%%%%%%%%%%%%%%%%%%%%%%%%%%%%%%%%%%

Profiling measures where a program spends computational resources.

A profiler may reveal:

\[
\boxed{
\text{runtime by function},
}
\]

\[
\boxed{
\text{memory allocation},
}
\]

\[
\boxed{
\text{communication time},
}
\]

or

\[
\boxed{
\text{hardware utilization}.
}
\]

Optimization should target measured bottlenecks rather than guessed ones.

%%%%%%%%%%%%%%%%%%%%%%%%%%%%%%%%%%%%%%%%%%%%%%%%%%%%%%%%%%%%
\subsection{Benchmarking}
\label{subsec:scientific-benchmarking}
%%%%%%%%%%%%%%%%%%%%%%%%%%%%%%%%%%%%%%%%%%%%%%%%%%%%%%%%%%%%

Benchmarking compares performance under controlled conditions.

Useful quantities include:

\[
\boxed{
\text{execution time},
}
\]

\[
\boxed{
\text{throughput},
}
\]

\[
\boxed{
\text{memory use},
}
\]

and

\[
\boxed{
\text{iterations to convergence}.
}
\]

A benchmark should specify the problem size, hardware, algorithm,
precision, tolerances, and implementation.

%%%%%%%%%%%%%%%%%%%%%%%%%%%%%%%%%%%%%%%%%%%%%%%%%%%%%%%%%%%%
\subsection{Warm-Up Effects}
\label{subsec:benchmark-warmup}
%%%%%%%%%%%%%%%%%%%%%%%%%%%%%%%%%%%%%%%%%%%%%%%%%%%%%%%%%%%%

The first execution of a computation may include additional costs such as:

\[
\boxed{
\text{memory initialization},
}
\]

\[
\boxed{
\text{cache population},
}
\]

or

\[
\boxed{
\text{runtime compilation}.
}
\]

Repeated measurements are often preferable to relying on one timing.

%%%%%%%%%%%%%%%%%%%%%%%%%%%%%%%%%%%%%%%%%%%%%%%%%%%%%%%%%%%%
\subsection{Parallel Computing}
\label{subsec:parallel-computing}
%%%%%%%%%%%%%%%%%%%%%%%%%%%%%%%%%%%%%%%%%%%%%%%%%%%%%%%%%%%%

Parallel computing performs multiple computations simultaneously.

The goal is to reduce runtime or solve larger problems by using several
processing units.

Parallelism may occur across:

\[
\boxed{
\text{cores},
}
\]

\[
\boxed{
\text{processors},
}
\]

\[
\boxed{
\text{GPUs},
}
\]

or

\[
\boxed{
\text{networked computers}.
}
\]

%%%%%%%%%%%%%%%%%%%%%%%%%%%%%%%%%%%%%%%%%%%%%%%%%%%%%%%%%%%%
\subsection{Task Parallelism}
\label{subsec:task-parallelism}
%%%%%%%%%%%%%%%%%%%%%%%%%%%%%%%%%%%%%%%%%%%%%%%%%%%%%%%%%%%%

Task parallelism performs different computations simultaneously.

For example:

\[
\boxed{
\text{simulation A}
\parallel
\text{simulation B}
\parallel
\text{simulation C}.
}
\]

Monte Carlo replications and parameter sweeps are naturally task
parallel.

%%%%%%%%%%%%%%%%%%%%%%%%%%%%%%%%%%%%%%%%%%%%%%%%%%%%%%%%%%%%
\subsection{Data Parallelism}
\label{subsec:data-parallelism}
%%%%%%%%%%%%%%%%%%%%%%%%%%%%%%%%%%%%%%%%%%%%%%%%%%%%%%%%%%%%

Data parallelism applies the same operation to different portions of a
data set simultaneously.

Examples include:

\[
\boxed{
\text{vector addition},
}
\]

\[
\boxed{
\text{matrix operations},
}
\]

and

\[
\boxed{
\text{grid updates in PDE solvers}.
}
\]

%%%%%%%%%%%%%%%%%%%%%%%%%%%%%%%%%%%%%%%%%%%%%%%%%%%%%%%%%%%%
\subsection{Shared-Memory Computing}
\label{subsec:shared-memory}
%%%%%%%%%%%%%%%%%%%%%%%%%%%%%%%%%%%%%%%%%%%%%%%%%%%%%%%%%%%%

In shared-memory systems, several processing units access a common memory
space.

Threads may operate on different portions of the same arrays.

Advantages include relatively simple data sharing.

Challenges include:

\[
\boxed{
\text{race conditions},
}
\]

\[
\boxed{
\text{synchronization},
}
\]

and

\[
\boxed{
\text{memory contention}.
}
\]

%%%%%%%%%%%%%%%%%%%%%%%%%%%%%%%%%%%%%%%%%%%%%%%%%%%%%%%%%%%%
\subsection{Race Conditions}
\label{subsec:race-conditions}
%%%%%%%%%%%%%%%%%%%%%%%%%%%%%%%%%%%%%%%%%%%%%%%%%%%%%%%%%%%%

A race condition occurs when the result depends on the timing of
simultaneous operations.

For example, two threads may attempt to update the same memory location.

Without synchronization, the final result may be unpredictable.

Parallel numerical algorithms should avoid or correctly manage shared
writes.

%%%%%%%%%%%%%%%%%%%%%%%%%%%%%%%%%%%%%%%%%%%%%%%%%%%%%%%%%%%%
\subsection{Synchronization}
\label{subsec:parallel-synchronization}
%%%%%%%%%%%%%%%%%%%%%%%%%%%%%%%%%%%%%%%%%%%%%%%%%%%%%%%%%%%%

Synchronization coordinates parallel tasks.

Examples include:

\[
\boxed{
\text{barriers},
}
\]

\[
\boxed{
\text{locks},
}
\]

and

\[
\boxed{
\text{collective operations}.
}
\]

Synchronization can ensure correctness but may reduce parallel efficiency.

%%%%%%%%%%%%%%%%%%%%%%%%%%%%%%%%%%%%%%%%%%%%%%%%%%%%%%%%%%%%
\subsection{Distributed-Memory Computing}
\label{subsec:distributed-memory}
%%%%%%%%%%%%%%%%%%%%%%%%%%%%%%%%%%%%%%%%%%%%%%%%%%%%%%%%%%%%

In distributed-memory systems, each process has its own memory.

Processes exchange information explicitly.

Conceptually,

\[
\boxed{
\text{local computation}
+
\text{message passing}.
}
\]

Large clusters and supercomputers commonly use this model.

%%%%%%%%%%%%%%%%%%%%%%%%%%%%%%%%%%%%%%%%%%%%%%%%%%%%%%%%%%%%
\subsection{Message Passing}
\label{subsec:message-passing}
%%%%%%%%%%%%%%%%%%%%%%%%%%%%%%%%%%%%%%%%%%%%%%%%%%%%%%%%%%%%

Processes may send and receive data:

\[
\boxed{
P_i
\longrightarrow
P_j.
}
\]

Communication may involve:

\[
\boxed{
\text{point-to-point messages}
}
\]

or

\[
\boxed{
\text{collective communication}.
}
\]

Collective operations include reductions, broadcasts, and global
synchronization.

%%%%%%%%%%%%%%%%%%%%%%%%%%%%%%%%%%%%%%%%%%%%%%%%%%%%%%%%%%%%
\subsection{Domain Decomposition}
\label{subsec:scientific-domain-decomposition}
%%%%%%%%%%%%%%%%%%%%%%%%%%%%%%%%%%%%%%%%%%%%%%%%%%%%%%%%%%%%

A computational domain can be divided among processors.

For example,

\[
\boxed{
\Omega
=
\Omega_1
\cup
\Omega_2
\cup
\cdots
\cup
\Omega_p.
}
\]

Each processor updates its local subdomain.

Boundary information is exchanged with neighboring subdomains.

This is common in numerical PDEs.

%%%%%%%%%%%%%%%%%%%%%%%%%%%%%%%%%%%%%%%%%%%%%%%%%%%%%%%%%%%%
\subsection{Halo and Ghost Data}
\label{subsec:halo-ghost-data}
%%%%%%%%%%%%%%%%%%%%%%%%%%%%%%%%%%%%%%%%%%%%%%%%%%%%%%%%%%%%

A processor may require neighboring values owned by another processor.

Copies of these values are stored in halo or ghost regions.

At each communication step, neighboring processes exchange updated
boundary data.

This permits local stencil computations.

%%%%%%%%%%%%%%%%%%%%%%%%%%%%%%%%%%%%%%%%%%%%%%%%%%%%%%%%%%%%
\subsection{Speedup}
\label{subsec:parallel-speedup}
%%%%%%%%%%%%%%%%%%%%%%%%%%%%%%%%%%%%%%%%%%%%%%%%%%%%%%%%%%%%

If one processor requires time

\[
T_1
\]

and \(p\) processors require time

\[
T_p,
\]

the parallel speedup is

\[
\boxed{
S_p
=
\frac{
T_1
}{
T_p
}.
}
\]

Ideal linear speedup is

\[
\boxed{
S_p=p.
}
\]

Real applications rarely achieve perfect linear scaling.

%%%%%%%%%%%%%%%%%%%%%%%%%%%%%%%%%%%%%%%%%%%%%%%%%%%%%%%%%%%%
\subsection{Parallel Efficiency}
\label{subsec:parallel-efficiency}
%%%%%%%%%%%%%%%%%%%%%%%%%%%%%%%%%%%%%%%%%%%%%%%%%%%%%%%%%%%%

Parallel efficiency is

\[
\boxed{
E_p
=
\frac{
S_p
}{
p
}
=
\frac{
T_1
}{
pT_p
}.
}
\]

Ideal efficiency is

\[
\boxed{
E_p=1.
}
\]

Lower efficiency indicates that additional processors are producing
diminishing benefit.

%%%%%%%%%%%%%%%%%%%%%%%%%%%%%%%%%%%%%%%%%%%%%%%%%%%%%%%%%%%%
\subsection{Worked Example: Speedup and Efficiency}
\label{subsec:worked-speedup-efficiency}
%%%%%%%%%%%%%%%%%%%%%%%%%%%%%%%%%%%%%%%%%%%%%%%%%%%%%%%%%%%%

\begin{workedexamplebox}{Parallel Performance}

Suppose a simulation takes

\[
T_1=120
\]

seconds on one processor and

\[
T_8=20
\]

seconds on eight processors.

The speedup is

\[
S_8
=
\frac{120}{20}
=
6.
\]

The parallel efficiency is

\[
E_8
=
\frac68
=
0.75.
\]

\begin{answerbox}
\[
\boxed{
S_8=6,
\qquad
E_8=75\%.
}
\]
\end{answerbox}

\end{workedexamplebox}

%%%%%%%%%%%%%%%%%%%%%%%%%%%%%%%%%%%%%%%%%%%%%%%%%%%%%%%%%%%%
\subsection{Amdahl's Law}
\label{subsec:amdahls-law}
%%%%%%%%%%%%%%%%%%%%%%%%%%%%%%%%%%%%%%%%%%%%%%%%%%%%%%%%%%%%

Suppose fraction \(s\) of a computation is inherently serial.

The remaining fraction

\[
1-s
\]

can be perfectly parallelized.

Then the ideal speedup on \(p\) processors is

\[
\boxed{
S_p
=
\frac{
1
}{
s+\frac{1-s}{p}
}.
}
\]

As

\[
p\to\infty,
\]

\[
\boxed{
S_{\max}
=
\frac1s.
}
\]

Thus even a small serial fraction can limit total speedup.

%%%%%%%%%%%%%%%%%%%%%%%%%%%%%%%%%%%%%%%%%%%%%%%%%%%%%%%%%%%%
\subsection{Worked Example: Amdahl's Law}
\label{subsec:worked-amdahl}
%%%%%%%%%%%%%%%%%%%%%%%%%%%%%%%%%%%%%%%%%%%%%%%%%%%%%%%%%%%%

\begin{workedexamplebox}{Maximum Parallel Speedup}

Suppose

\[
s=0.10.
\]

Then with infinitely many processors,

\[
S_{\max}
=
\frac1{0.10}
=
10.
\]

\begin{answerbox}
\[
\boxed{
S_{\max}=10.
}
\]
\end{answerbox}

No amount of additional parallel hardware can exceed this idealized limit
unless the serial portion is reduced.

\end{workedexamplebox}

%%%%%%%%%%%%%%%%%%%%%%%%%%%%%%%%%%%%%%%%%%%%%%%%%%%%%%%%%%%%
\subsection{Strong Scaling}
\label{subsec:strong-scaling}
%%%%%%%%%%%%%%%%%%%%%%%%%%%%%%%%%%%%%%%%%%%%%%%%%%%%%%%%%%%%

Strong scaling keeps total problem size fixed while increasing processor
count.

The question is:

\[
\boxed{
\text{How much faster can the same problem be solved?}
}
\]

Eventually, communication and synchronization dominate.

%%%%%%%%%%%%%%%%%%%%%%%%%%%%%%%%%%%%%%%%%%%%%%%%%%%%%%%%%%%%
\subsection{Weak Scaling}
\label{subsec:weak-scaling}
%%%%%%%%%%%%%%%%%%%%%%%%%%%%%%%%%%%%%%%%%%%%%%%%%%%%%%%%%%%%

Weak scaling increases problem size proportionally with processor count.

The amount of work per processor remains approximately constant.

The question is:

\[
\boxed{
\text{Can a larger problem be solved in approximately the same time?}
}
\]

%%%%%%%%%%%%%%%%%%%%%%%%%%%%%%%%%%%%%%%%%%%%%%%%%%%%%%%%%%%%
\subsection{Load Balancing}
\label{subsec:load-balancing}
%%%%%%%%%%%%%%%%%%%%%%%%%%%%%%%%%%%%%%%%%%%%%%%%%%%%%%%%%%%%

Parallel efficiency decreases if some processors perform much more work
than others.

Ideally,

\[
\boxed{
W_1
\approx
W_2
\approx
\cdots
\approx
W_p.
}
\]

Poor load balancing causes processors to wait for the slowest worker.

%%%%%%%%%%%%%%%%%%%%%%%%%%%%%%%%%%%%%%%%%%%%%%%%%%%%%%%%%%%%
\subsection{Communication Overhead}
\label{subsec:communication-overhead}
%%%%%%%%%%%%%%%%%%%%%%%%%%%%%%%%%%%%%%%%%%%%%%%%%%%%%%%%%%%%

Parallel computation often requires communication.

Total runtime may be modeled schematically as

\[
\boxed{
T
=
T_{\mathrm{compute}}
+
T_{\mathrm{communication}}
+
T_{\mathrm{synchronization}}.
}
\]

Increasing processor count may reduce computational work per processor
while increasing communication cost.

%%%%%%%%%%%%%%%%%%%%%%%%%%%%%%%%%%%%%%%%%%%%%%%%%%%%%%%%%%%%
\subsection{Latency and Bandwidth}
\label{subsec:latency-bandwidth}
%%%%%%%%%%%%%%%%%%%%%%%%%%%%%%%%%%%%%%%%%%%%%%%%%%%%%%%%%%%%

Communication cost is often described by:

\[
\boxed{
\text{latency}
=
\text{startup cost per message},
}
\]

and

\[
\boxed{
\text{bandwidth}
=
\text{data transfer rate}.
}
\]

Many small messages can be inefficient because latency is paid repeatedly.

%%%%%%%%%%%%%%%%%%%%%%%%%%%%%%%%%%%%%%%%%%%%%%%%%%%%%%%%%%%%
\subsection{Granularity}
\label{subsec:parallel-granularity}
%%%%%%%%%%%%%%%%%%%%%%%%%%%%%%%%%%%%%%%%%%%%%%%%%%%%%%%%%%%%

Granularity describes the amount of computation performed between
communication or synchronization events.

Fine-grained parallelism uses many small tasks.

Coarse-grained parallelism uses larger tasks.

Too-fine granularity may cause overhead to dominate useful computation.

%%%%%%%%%%%%%%%%%%%%%%%%%%%%%%%%%%%%%%%%%%%%%%%%%%%%%%%%%%%%
\subsection{Embarrassingly Parallel Problems}
\label{subsec:embarrassingly-parallel}
%%%%%%%%%%%%%%%%%%%%%%%%%%%%%%%%%%%%%%%%%%%%%%%%%%%%%%%%%%%%

A problem is called embarrassingly parallel when tasks require little or no
communication.

Examples include:

\[
\boxed{
\text{independent Monte Carlo paths},
}
\]

\[
\boxed{
\text{parameter sweeps},
}
\]

and

\[
\boxed{
\text{independent optimization starts}.
}
\]

Such problems can achieve excellent parallel scaling.

%%%%%%%%%%%%%%%%%%%%%%%%%%%%%%%%%%%%%%%%%%%%%%%%%%%%%%%%%%%%
\subsection{Graphics Processing Units}
\label{subsec:gpu-computing}
%%%%%%%%%%%%%%%%%%%%%%%%%%%%%%%%%%%%%%%%%%%%%%%%%%%%%%%%%%%%

Graphics Processing Units, abbreviated GPUs, contain many computational
units designed for highly parallel workloads.

They are effective when many similar operations can be performed
simultaneously.

Examples include:

\[
\boxed{
\text{dense matrix operations},
}
\]

\[
\boxed{
\text{large vector operations},
}
\]

and

\[
\boxed{
\text{independent simulation kernels}.
}
\]

%%%%%%%%%%%%%%%%%%%%%%%%%%%%%%%%%%%%%%%%%%%%%%%%%%%%%%%%%%%%
\subsection{CPU Versus GPU}
\label{subsec:cpu-versus-gpu}
%%%%%%%%%%%%%%%%%%%%%%%%%%%%%%%%%%%%%%%%%%%%%%%%%%%%%%%%%%%%

A CPU typically contains fewer powerful general-purpose cores.

A GPU contains many simpler processing units optimized for throughput.

Thus:

\[
\boxed{
\text{CPU}
\rightarrow
\text{complex control and general tasks},
}
\]

while

\[
\boxed{
\text{GPU}
\rightarrow
\text{massively data-parallel work}.
}
\]

The distinction is architectural rather than absolute.

%%%%%%%%%%%%%%%%%%%%%%%%%%%%%%%%%%%%%%%%%%%%%%%%%%%%%%%%%%%%
\subsection{Accelerator Data Transfer}
\label{subsec:accelerator-data-transfer}
%%%%%%%%%%%%%%%%%%%%%%%%%%%%%%%%%%%%%%%%%%%%%%%%%%%%%%%%%%%%

Moving data between processor memory spaces can be expensive.

A computation may lose the advantage of an accelerator if arrays are
repeatedly transferred.

Efficient accelerator algorithms attempt to:

\[
\boxed{
\text{move data infrequently}
}
\]

and

\[
\boxed{
\text{perform substantial work while data remain resident}.
}
\]

%%%%%%%%%%%%%%%%%%%%%%%%%%%%%%%%%%%%%%%%%%%%%%%%%%%%%%%%%%%%
\subsection{Precision}
\label{subsec:scientific-computing-precision}
%%%%%%%%%%%%%%%%%%%%%%%%%%%%%%%%%%%%%%%%%%%%%%%%%%%%%%%%%%%%

Scientific computations may use several numerical precisions.

Common choices include approximately:

\[
\boxed{
\text{single precision},
}
\]

and

\[
\boxed{
\text{double precision}.
}
\]

Higher or arbitrary precision can be used when ordinary floating-point
precision is insufficient.

%%%%%%%%%%%%%%%%%%%%%%%%%%%%%%%%%%%%%%%%%%%%%%%%%%%%%%%%%%%%
\subsection{Single and Double Precision}
\label{subsec:single-double-precision}
%%%%%%%%%%%%%%%%%%%%%%%%%%%%%%%%%%%%%%%%%%%%%%%%%%%%%%%%%%%%

Single precision uses fewer bits and therefore:

\[
\boxed{
\text{less memory}
}
\]

and often

\[
\boxed{
\text{greater computational throughput}.
}
\]

Double precision provides substantially greater numerical accuracy and
dynamic range.

The appropriate choice depends on conditioning and required accuracy.

%%%%%%%%%%%%%%%%%%%%%%%%%%%%%%%%%%%%%%%%%%%%%%%%%%%%%%%%%%%%
\subsection{Mixed Precision}
\label{subsec:scientific-mixed-precision}
%%%%%%%%%%%%%%%%%%%%%%%%%%%%%%%%%%%%%%%%%%%%%%%%%%%%%%%%%%%%

Mixed-precision algorithms use different precisions for different parts of
a computation.

For example:

\[
\boxed{
\text{low-precision factorization}
}
\]

may be combined with

\[
\boxed{
\text{high-precision residual correction}.
}
\]

Iterative refinement is an important example.

%%%%%%%%%%%%%%%%%%%%%%%%%%%%%%%%%%%%%%%%%%%%%%%%%%%%%%%%%%%%
\subsection{Why Lower Precision Can Be Faster}
\label{subsec:why-lower-precision-faster}
%%%%%%%%%%%%%%%%%%%%%%%%%%%%%%%%%%%%%%%%%%%%%%%%%%%%%%%%%%%%

Lower-precision values require fewer bits.

This may reduce:

\[
\boxed{
\text{memory use},
}
\]

\[
\boxed{
\text{memory bandwidth},
}
\]

and sometimes

\[
\boxed{
\text{arithmetic cost}.
}
\]

The benefit is useful only if the resulting numerical accuracy remains
acceptable.

%%%%%%%%%%%%%%%%%%%%%%%%%%%%%%%%%%%%%%%%%%%%%%%%%%%%%%%%%%%%
\subsection{Arbitrary Precision}
\label{subsec:arbitrary-precision}
%%%%%%%%%%%%%%%%%%%%%%%%%%%%%%%%%%%%%%%%%%%%%%%%%%%%%%%%%%%%

Arbitrary-precision arithmetic allows the number of stored digits to be
increased beyond standard hardware precision.

It is useful for:

\[
\boxed{
\text{ill-conditioned problems},
}
\]

\[
\boxed{
\text{high-accuracy constants},
}
\]

or

\[
\boxed{
\text{verification of numerical conjectures}.
}
\]

The computational cost increases significantly with precision.

%%%%%%%%%%%%%%%%%%%%%%%%%%%%%%%%%%%%%%%%%%%%%%%%%%%%%%%%%%%%
\subsection{Floating-Point Reduction Order}
\label{subsec:floating-point-reduction-order}
%%%%%%%%%%%%%%%%%%%%%%%%%%%%%%%%%%%%%%%%%%%%%%%%%%%%%%%%%%%%

Floating-point addition is not exactly associative:

\[
\boxed{
(a+b)+c
\neq
a+(b+c)
}
\]

in general.

Parallel reductions may sum values in different orders.

Thus two parallel runs can produce slightly different final digits even
when the mathematical computation is equivalent.

%%%%%%%%%%%%%%%%%%%%%%%%%%%%%%%%%%%%%%%%%%%%%%%%%%%%%%%%%%%%
\subsection{Pairwise Summation}
\label{subsec:pairwise-summation}
%%%%%%%%%%%%%%%%%%%%%%%%%%%%%%%%%%%%%%%%%%%%%%%%%%%%%%%%%%%%

Instead of summing sequentially,

\[
x_1+x_2+\cdots+x_n,
\]

pairwise summation recursively sums groups of similar size.

This can reduce accumulated floating-point error and also maps naturally
onto parallel reduction trees.

%%%%%%%%%%%%%%%%%%%%%%%%%%%%%%%%%%%%%%%%%%%%%%%%%%%%%%%%%%%%
\subsection{Compensated Summation}
\label{subsec:compensated-summation}
%%%%%%%%%%%%%%%%%%%%%%%%%%%%%%%%%%%%%%%%%%%%%%%%%%%%%%%%%%%%

Compensated summation attempts to track rounding error lost during repeated
addition.

Conceptually, one maintains both:

\[
\boxed{
\text{running sum}
}
\]

and

\[
\boxed{
\text{small correction}.
}
\]

This can greatly improve accuracy when summing many floating-point values.

%%%%%%%%%%%%%%%%%%%%%%%%%%%%%%%%%%%%%%%%%%%%%%%%%%%%%%%%%%%%
\subsection{Deterministic and Nondeterministic Computation}
\label{subsec:deterministic-nondeterministic}
%%%%%%%%%%%%%%%%%%%%%%%%%%%%%%%%%%%%%%%%%%%%%%%%%%%%%%%%%%%%

A deterministic program produces the same mathematical execution path from
the same inputs.

Parallel scheduling, randomized algorithms, asynchronous updates, and
hardware-dependent reduction orders may introduce nondeterministic
execution.

This does not automatically imply mathematical invalidity.

But reproducibility must be interpreted carefully.

%%%%%%%%%%%%%%%%%%%%%%%%%%%%%%%%%%%%%%%%%%%%%%%%%%%%%%%%%%%%
\subsection{Bitwise Reproducibility}
\label{subsec:bitwise-reproducibility}
%%%%%%%%%%%%%%%%%%%%%%%%%%%%%%%%%%%%%%%%%%%%%%%%%%%%%%%%%%%%

Bitwise reproducibility means repeated runs produce exactly identical
binary outputs.

This is stronger than numerical reproducibility.

Two correct calculations may differ in low-order floating-point digits
while agreeing to all scientifically meaningful precision.

%%%%%%%%%%%%%%%%%%%%%%%%%%%%%%%%%%%%%%%%%%%%%%%%%%%%%%%%%%%%
\subsection{Numerical Reproducibility}
\label{subsec:numerical-reproducibility}
%%%%%%%%%%%%%%%%%%%%%%%%%%%%%%%%%%%%%%%%%%%%%%%%%%%%%%%%%%%%

Numerical reproducibility means independently repeated computations agree
within stated numerical tolerances.

This often provides a more meaningful scientific standard than exact
bitwise equality.

The acceptable tolerance should be connected to:

\[
\boxed{
\text{conditioning},
}
\]

\[
\boxed{
\text{discretization error},
}
\]

and

\[
\boxed{
\text{scientific precision requirements}.
}
\]

%%%%%%%%%%%%%%%%%%%%%%%%%%%%%%%%%%%%%%%%%%%%%%%%%%%%%%%%%%%%
\subsection{Reproducible Computational Experiments}
\label{subsec:reproducible-computational-experiments}
%%%%%%%%%%%%%%%%%%%%%%%%%%%%%%%%%%%%%%%%%%%%%%%%%%%%%%%%%%%%

A computational experiment should record:

\[
\boxed{
\text{model equations},
}
\]

\[
\boxed{
\text{input data},
}
\]

\[
\boxed{
\text{parameter values},
}
\]

\[
\boxed{
\text{software versions},
}
\]

\[
\boxed{
\text{numerical tolerances},
}
\]

\[
\boxed{
\text{random seeds},
}
\]

and

\[
\boxed{
\text{hardware information when relevant}.
}
\]

%%%%%%%%%%%%%%%%%%%%%%%%%%%%%%%%%%%%%%%%%%%%%%%%%%%%%%%%%%%%
\subsection{Data Provenance}
\label{subsec:data-provenance}
%%%%%%%%%%%%%%%%%%%%%%%%%%%%%%%%%%%%%%%%%%%%%%%%%%%%%%%%%%%%

Data provenance records where data came from and how they were processed.

A derived data set should ideally be traceable through

\[
\boxed{
\text{raw input}
\rightarrow
\text{transformations}
\rightarrow
\text{analysis}
\rightarrow
\text{final result}.
}
\]

This is especially important in long computational pipelines.

%%%%%%%%%%%%%%%%%%%%%%%%%%%%%%%%%%%%%%%%%%%%%%%%%%%%%%%%%%%%
\subsection{Configuration Files}
\label{subsec:scientific-configuration-files}
%%%%%%%%%%%%%%%%%%%%%%%%%%%%%%%%%%%%%%%%%%%%%%%%%%%%%%%%%%%%

Large simulations often separate parameters from source code.

A configuration may contain:

\[
\boxed{
\text{mesh size},
}
\]

\[
\boxed{
\text{time step},
}
\]

\[
\boxed{
\text{solver tolerance},
}
\]

\[
\boxed{
\text{physical parameters},
}
\]

and

\[
\boxed{
\text{output options}.
}
\]

This improves experiment management and reproducibility.

%%%%%%%%%%%%%%%%%%%%%%%%%%%%%%%%%%%%%%%%%%%%%%%%%%%%%%%%%%%%
\subsection{Logging}
\label{subsec:scientific-logging}
%%%%%%%%%%%%%%%%%%%%%%%%%%%%%%%%%%%%%%%%%%%%%%%%%%%%%%%%%%%%

A simulation log may record:

\[
\boxed{
\text{start and end times},
}
\]

\[
\boxed{
\text{iteration counts},
}
\]

\[
\boxed{
\text{residual histories},
}
\]

\[
\boxed{
\text{warnings},
}
\]

and

\[
\boxed{
\text{parameter settings}.
}
\]

Logs help diagnose both numerical and software failures.

%%%%%%%%%%%%%%%%%%%%%%%%%%%%%%%%%%%%%%%%%%%%%%%%%%%%%%%%%%%%
\subsection{Checkpointing}
\label{subsec:checkpointing}
%%%%%%%%%%%%%%%%%%%%%%%%%%%%%%%%%%%%%%%%%%%%%%%%%%%%%%%%%%%%

Long computations should periodically save enough state to restart.

A checkpoint may contain:

\[
\boxed{
\text{current solution},
}
\]

\[
\boxed{
\text{simulation time},
}
\]

\[
\boxed{
\text{solver state},
}
\]

and possibly

\[
\boxed{
\text{random generator state}.
}
\]

If the program stops unexpectedly, computation can resume from the latest
checkpoint.

%%%%%%%%%%%%%%%%%%%%%%%%%%%%%%%%%%%%%%%%%%%%%%%%%%%%%%%%%%%%
\subsection{Checkpoint Interval}
\label{subsec:checkpoint-interval}
%%%%%%%%%%%%%%%%%%%%%%%%%%%%%%%%%%%%%%%%%%%%%%%%%%%%%%%%%%%%

Saving checkpoints too frequently creates excessive I/O cost.

Saving them too rarely risks losing substantial computation after a
failure.

Thus checkpoint frequency balances:

\[
\boxed{
\text{storage overhead}
}
\]

against

\[
\boxed{
\text{recomputation risk}.
}
\]

%%%%%%%%%%%%%%%%%%%%%%%%%%%%%%%%%%%%%%%%%%%%%%%%%%%%%%%%%%%%
\subsection{Parameter Sweeps}
\label{subsec:parameter-sweeps}
%%%%%%%%%%%%%%%%%%%%%%%%%%%%%%%%%%%%%%%%%%%%%%%%%%%%%%%%%%%%

Suppose a model depends on parameter

\[
\theta.
\]

A parameter sweep evaluates

\[
\boxed{
M(\theta_1),
M(\theta_2),
\ldots,
M(\theta_N).
}
\]

Parameter sweeps are useful for:

\[
\boxed{
\text{sensitivity analysis},
}
\]

\[
\boxed{
\text{phase diagrams},
}
\]

and

\[
\boxed{
\text{scenario comparison}.
}
\]

Independent parameter combinations are naturally parallel.

%%%%%%%%%%%%%%%%%%%%%%%%%%%%%%%%%%%%%%%%%%%%%%%%%%%%%%%%%%%%
\subsection{Ensemble Simulation}
\label{subsec:ensemble-simulation}
%%%%%%%%%%%%%%%%%%%%%%%%%%%%%%%%%%%%%%%%%%%%%%%%%%%%%%%%%%%%

An ensemble consists of many simulations with different:

\[
\boxed{
\text{initial conditions},
}
\]

\[
\boxed{
\text{parameters},
}
\]

or

\[
\boxed{
\text{random inputs}.
}
\]

Ensemble results provide information about variability and uncertainty.

Monte Carlo simulation is an important form of ensemble computation.

%%%%%%%%%%%%%%%%%%%%%%%%%%%%%%%%%%%%%%%%%%%%%%%%%%%%%%%%%%%%
\subsection{Batch Processing}
\label{subsec:batch-processing}
%%%%%%%%%%%%%%%%%%%%%%%%%%%%%%%%%%%%%%%%%%%%%%%%%%%%%%%%%%%%

Large computational experiments may submit many independent jobs to a
shared computing system.

Each job processes a particular:

\[
\boxed{
\text{parameter case},
}
\]

\[
\boxed{
\text{sample},
}
\]

or

\[
\boxed{
\text{simulation segment}.
}
\]

Batch systems schedule these jobs based on available resources.

%%%%%%%%%%%%%%%%%%%%%%%%%%%%%%%%%%%%%%%%%%%%%%%%%%%%%%%%%%%%
\subsection{Computational Pipelines}
\label{subsec:computational-pipelines}
%%%%%%%%%%%%%%%%%%%%%%%%%%%%%%%%%%%%%%%%%%%%%%%%%%%%%%%%%%%%

A scientific workflow often contains several stages:

\[
\boxed{
\text{input}
\rightarrow
\text{preprocessing}
\rightarrow
\text{simulation}
\rightarrow
\text{analysis}
\rightarrow
\text{visualization}.
}
\]

Automating these stages reduces manual inconsistency and improves
reproducibility.

%%%%%%%%%%%%%%%%%%%%%%%%%%%%%%%%%%%%%%%%%%%%%%%%%%%%%%%%%%%%
\subsection{Workflow Dependencies}
\label{subsec:workflow-dependencies}
%%%%%%%%%%%%%%%%%%%%%%%%%%%%%%%%%%%%%%%%%%%%%%%%%%%%%%%%%%%%

Some computational tasks cannot begin until others finish.

This forms a dependency graph.

For example,

\[
\boxed{
\text{mesh generation}
\rightarrow
\text{PDE solve}
\rightarrow
\text{postprocessing}.
}
\]

Independent branches of the graph can execute in parallel.

%%%%%%%%%%%%%%%%%%%%%%%%%%%%%%%%%%%%%%%%%%%%%%%%%%%%%%%%%%%%
\subsection{Verification of Scientific Software}
\label{subsec:scientific-software-verification}
%%%%%%%%%%%%%%%%%%%%%%%%%%%%%%%%%%%%%%%%%%%%%%%%%%%%%%%%%%%%

Software verification asks whether the implemented code correctly performs
the intended numerical method.

Verification techniques include:

\[
\boxed{
\text{unit tests},
}
\]

\[
\boxed{
\text{analytic test cases},
}
\]

\[
\boxed{
\text{manufactured solutions},
}
\]

\[
\boxed{
\text{convergence studies},
}
\]

and

\[
\boxed{
\text{cross-code comparisons}.
}
\]

%%%%%%%%%%%%%%%%%%%%%%%%%%%%%%%%%%%%%%%%%%%%%%%%%%%%%%%%%%%%
\subsection{Unit Testing}
\label{subsec:scientific-unit-testing}
%%%%%%%%%%%%%%%%%%%%%%%%%%%%%%%%%%%%%%%%%%%%%%%%%%%%%%%%%%%%

A unit test checks one small component of a program.

Examples include testing:

\[
\boxed{
\text{a derivative routine},
}
\]

\[
\boxed{
\text{a matrix assembly function},
}
\]

or

\[
\boxed{
\text{a coordinate transformation}.
}
\]

Small tests make software failures easier to isolate.

%%%%%%%%%%%%%%%%%%%%%%%%%%%%%%%%%%%%%%%%%%%%%%%%%%%%%%%%%%%%
\subsection{Regression Testing}
\label{subsec:regression-testing}
%%%%%%%%%%%%%%%%%%%%%%%%%%%%%%%%%%%%%%%%%%%%%%%%%%%%%%%%%%%%

A regression test compares current software output with previously
validated output.

Its purpose is to detect unintended changes after modifying the code.

Because floating-point outputs may vary slightly, numerical regression
tests should often use tolerances rather than exact string equality.

%%%%%%%%%%%%%%%%%%%%%%%%%%%%%%%%%%%%%%%%%%%%%%%%%%%%%%%%%%%%
\subsection{Convergence Testing}
\label{subsec:scientific-convergence-testing}
%%%%%%%%%%%%%%%%%%%%%%%%%%%%%%%%%%%%%%%%%%%%%%%%%%%%%%%%%%%%

Suppose a numerical method is theoretically order \(p\).

A convergence test checks whether

\[
\boxed{
E(h)
\approx
Ch^p.
}
\]

Then

\[
\boxed{
\frac{
E(h)
}{
E(h/2)
}
\approx
2^p.
}
\]

Failure to observe the expected rate can reveal implementation or modeling
errors.

%%%%%%%%%%%%%%%%%%%%%%%%%%%%%%%%%%%%%%%%%%%%%%%%%%%%%%%%%%%%
\subsection{Residual Testing}
\label{subsec:scientific-residual-testing}
%%%%%%%%%%%%%%%%%%%%%%%%%%%%%%%%%%%%%%%%%%%%%%%%%%%%%%%%%%%%

For a computed solution \(u_h\), evaluate the discrete residual

\[
\boxed{
r_h
=
F_h(u_h).
}
\]

A large residual indicates the algebraic equations have not been solved
accurately.

A small residual does not by itself guarantee small physical or
discretization error.

%%%%%%%%%%%%%%%%%%%%%%%%%%%%%%%%%%%%%%%%%%%%%%%%%%%%%%%%%%%%
\subsection{Conservation Checks}
\label{subsec:scientific-conservation-checks}
%%%%%%%%%%%%%%%%%%%%%%%%%%%%%%%%%%%%%%%%%%%%%%%%%%%%%%%%%%%%

If the mathematical model conserves a quantity

\[
I(u),
\]

monitor

\[
\boxed{
I(u^n)-I(u^0).
}
\]

Examples include:

\[
\boxed{
\text{mass},
}
\]

\[
\boxed{
\text{charge},
}
\]

\[
\boxed{
\text{energy},
}
\]

or

\[
\boxed{
\text{population}.
}
\]

Unexpected drift can reveal numerical or software errors.

%%%%%%%%%%%%%%%%%%%%%%%%%%%%%%%%%%%%%%%%%%%%%%%%%%%%%%%%%%%%
\subsection{Dimensional Analysis as a Software Check}
\label{subsec:dimensional-analysis-software-check}
%%%%%%%%%%%%%%%%%%%%%%%%%%%%%%%%%%%%%%%%%%%%%%%%%%%%%%%%%%%%

Physical equations must be dimensionally consistent.

For example, if

\[
x
=
x_0
+
vt,
\]

then

\[
[v][t]
=
[x].
\]

Unit inconsistencies can reveal input and implementation mistakes before
expensive simulations are run.

%%%%%%%%%%%%%%%%%%%%%%%%%%%%%%%%%%%%%%%%%%%%%%%%%%%%%%%%%%%%
\subsection{Nondimensionalization}
\label{subsec:scientific-nondimensionalization}
%%%%%%%%%%%%%%%%%%%%%%%%%%%%%%%%%%%%%%%%%%%%%%%%%%%%%%%%%%%%

Introduce characteristic scales such as

\[
x=L\widehat{x},
\qquad
t=T\widehat{t},
\qquad
u=U\widehat{u}.
\]

The resulting equations may depend on dimensionless parameters.

Nondimensionalization can improve:

\[
\boxed{
\text{physical interpretation},
}
\]

and

\[
\boxed{
\text{numerical scaling}.
}
\]

%%%%%%%%%%%%%%%%%%%%%%%%%%%%%%%%%%%%%%%%%%%%%%%%%%%%%%%%%%%%
\subsection{Example: Dimensionless ODE}
\label{subsec:worked-dimensionless-ode}
%%%%%%%%%%%%%%%%%%%%%%%%%%%%%%%%%%%%%%%%%%%%%%%%%%%%%%%%%%%%

\begin{workedexamplebox}{Scaling Exponential Decay}

Suppose

\[
\frac{dy}{dt}
=
-ky.
\]

Introduce

\[
\tau=kt.
\]

Then

\[
\frac{dy}{dt}
=
k
\frac{dy}{d\tau}.
\]

Therefore,

\[
k
\frac{dy}{d\tau}
=
-ky.
\]

Canceling \(k\),

\[
\begin{answerbox}
\[
\boxed{
\frac{dy}{d\tau}
=
-y.
}
\]
\end{answerbox}

The physical parameter has been absorbed into the time scale.

\end{workedexamplebox}

%%%%%%%%%%%%%%%%%%%%%%%%%%%%%%%%%%%%%%%%%%%%%%%%%%%%%%%%%%%%
\subsection{Sensitivity Studies}
\label{subsec:scientific-sensitivity-studies}
%%%%%%%%%%%%%%%%%%%%%%%%%%%%%%%%%%%%%%%%%%%%%%%%%%%%%%%%%%%%

A sensitivity study varies inputs and observes the resulting output.

For parameter \(\theta\),

\[
\boxed{
S
=
\frac{
\partial Q
}{
\partial\theta
}
}
\]

measures local sensitivity of output \(Q\).

Finite differences may approximate:

\[
\boxed{
\frac{
Q(\theta+h)-Q(\theta-h)
}{
2h
}.
}
\]

%%%%%%%%%%%%%%%%%%%%%%%%%%%%%%%%%%%%%%%%%%%%%%%%%%%%%%%%%%%%
\subsection{Uncertainty Versus Sensitivity}
\label{subsec:uncertainty-versus-sensitivity}
%%%%%%%%%%%%%%%%%%%%%%%%%%%%%%%%%%%%%%%%%%%%%%%%%%%%%%%%%%%%

Sensitivity asks how strongly an output responds to a parameter.

Uncertainty quantification asks how uncertain inputs produce uncertain
outputs.

A parameter can have high sensitivity but little uncertainty, or low
sensitivity but large uncertainty.

These concepts should not be confused.

%%%%%%%%%%%%%%%%%%%%%%%%%%%%%%%%%%%%%%%%%%%%%%%%%%%%%%%%%%%%
\subsection{Scientific Visualization}
\label{subsec:scientific-visualization}
%%%%%%%%%%%%%%%%%%%%%%%%%%%%%%%%%%%%%%%%%%%%%%%%%%%%%%%%%%%%

Visualization helps interpret large numerical data sets.

Common visualizations include:

\[
\boxed{
\text{line plots},
}
\]

\[
\boxed{
\text{contours},
}
\]

\[
\boxed{
\text{vector fields},
}
\]

\[
\boxed{
\text{surface plots},
}
\]

and

\[
\boxed{
\text{animations}.
}
\]

Visualization can reveal patterns that are difficult to detect in raw
arrays.

%%%%%%%%%%%%%%%%%%%%%%%%%%%%%%%%%%%%%%%%%%%%%%%%%%%%%%%%%%%%
\subsection{Visualization Is Not Validation}
\label{subsec:visualization-not-validation}
%%%%%%%%%%%%%%%%%%%%%%%%%%%%%%%%%%%%%%%%%%%%%%%%%%%%%%%%%%%%

A visually smooth solution may still be:

\[
\boxed{
\text{inaccurate},
}
\]

\[
\boxed{
\text{unstable},
}
\]

or

\[
\boxed{
\text{physically incorrect}.
}
\]

Plots should be combined with quantitative diagnostics.

%%%%%%%%%%%%%%%%%%%%%%%%%%%%%%%%%%%%%%%%%%%%%%%%%%%%%%%%%%%%
\subsection{Derived Quantities}
\label{subsec:derived-quantities}
%%%%%%%%%%%%%%%%%%%%%%%%%%%%%%%%%%%%%%%%%%%%%%%%%%%%%%%%%%%%

Scientific interest often lies not in the full solution field but in
derived quantities.

Examples include:

\[
\boxed{
\text{maximum temperature},
}
\]

\[
\boxed{
\text{total flux},
}
\]

\[
\boxed{
\text{mean displacement},
}
\]

or

\[
\boxed{
\text{peak infection}.
}
\]

Error analysis should consider these quantities of interest directly.

%%%%%%%%%%%%%%%%%%%%%%%%%%%%%%%%%%%%%%%%%%%%%%%%%%%%%%%%%%%%
\subsection{Reduction Operations}
\label{subsec:scientific-reductions}
%%%%%%%%%%%%%%%%%%%%%%%%%%%%%%%%%%%%%%%%%%%%%%%%%%%%%%%%%%%%

A reduction combines many values into one.

Examples include:

\[
\boxed{
\sum_i x_i,
}
\]

\[
\boxed{
\max_i x_i,
}
\]

and

\[
\boxed{
\|x\|_2.
}
\]

Parallel reductions require communication and may be synchronization
points.

%%%%%%%%%%%%%%%%%%%%%%%%%%%%%%%%%%%%%%%%%%%%%%%%%%%%%%%%%%%%
\subsection{Input and Output Costs}
\label{subsec:scientific-io}
%%%%%%%%%%%%%%%%%%%%%%%%%%%%%%%%%%%%%%%%%%%%%%%%%%%%%%%%%%%%

Large simulations can generate enormous amounts of data.

Writing every state to storage may dominate runtime.

Output strategies may instead save:

\[
\boxed{
\text{selected time steps},
}
\]

\[
\boxed{
\text{compressed fields},
}
\]

or

\[
\boxed{
\text{derived quantities}.
}
\]

%%%%%%%%%%%%%%%%%%%%%%%%%%%%%%%%%%%%%%%%%%%%%%%%%%%%%%%%%%%%
\subsection{In Situ Analysis}
\label{subsec:in-situ-analysis}
%%%%%%%%%%%%%%%%%%%%%%%%%%%%%%%%%%%%%%%%%%%%%%%%%%%%%%%%%%%%

In situ analysis processes simulation data while the simulation is still
running.

Instead of writing every full data set to disk, the program computes
summaries, statistics, or visualization information immediately.

This can reduce storage and I/O requirements.

%%%%%%%%%%%%%%%%%%%%%%%%%%%%%%%%%%%%%%%%%%%%%%%%%%%%%%%%%%%%
\subsection{Computational Storage}
\label{subsec:computational-storage}
%%%%%%%%%%%%%%%%%%%%%%%%%%%%%%%%%%%%%%%%%%%%%%%%%%%%%%%%%%%%

Scientific data may include:

\[
\boxed{
\text{raw inputs},
}
\]

\[
\boxed{
\text{checkpoints},
}
\]

\[
\boxed{
\text{solution fields},
}
\]

\[
\boxed{
\text{logs},
}
\]

and

\[
\boxed{
\text{derived outputs}.
}
\]

A storage strategy should preserve enough information to reproduce key
scientific conclusions without retaining unnecessary data.

%%%%%%%%%%%%%%%%%%%%%%%%%%%%%%%%%%%%%%%%%%%%%%%%%%%%%%%%%%%%
\subsection{High-Performance Computing}
\label{subsec:high-performance-computing}
%%%%%%%%%%%%%%%%%%%%%%%%%%%%%%%%%%%%%%%%%%%%%%%%%%%%%%%%%%%%

High-Performance Computing, abbreviated HPC, uses advanced computational
resources to solve problems requiring substantial:

\[
\boxed{
\text{processing power},
}
\]

\[
\boxed{
\text{memory},
}
\]

or

\[
\boxed{
\text{parallelism}.
}
\]

HPC includes multicore processors, GPUs, clusters, and supercomputers.

%%%%%%%%%%%%%%%%%%%%%%%%%%%%%%%%%%%%%%%%%%%%%%%%%%%%%%%%%%%%
\subsection{Scalability}
\label{subsec:scientific-scalability}
%%%%%%%%%%%%%%%%%%%%%%%%%%%%%%%%%%%%%%%%%%%%%%%%%%%%%%%%%%%%

A scalable algorithm continues to use additional computational resources
effectively as:

\[
\boxed{
\text{problem size increases}
}
\]

or

\[
\boxed{
\text{processor count increases}.
}
\]

Scalability depends on both the mathematical algorithm and its
implementation.

%%%%%%%%%%%%%%%%%%%%%%%%%%%%%%%%%%%%%%%%%%%%%%%%%%%%%%%%%%%%
\subsection{Algorithmic Scalability}
\label{subsec:algorithmic-scalability}
%%%%%%%%%%%%%%%%%%%%%%%%%%%%%%%%%%%%%%%%%%%%%%%%%%%%%%%%%%%%

Suppose one solver requires

\[
O(N^3)
\]

work while another requires

\[
O(N).
\]

For small \(N\), the difference may not matter.

For very large \(N\), asymptotic complexity dominates.

Thus better algorithms are often more important than faster hardware.

%%%%%%%%%%%%%%%%%%%%%%%%%%%%%%%%%%%%%%%%%%%%%%%%%%%%%%%%%%%%
\subsection{Parallel Scalability}
\label{subsec:parallel-scalability}
%%%%%%%%%%%%%%%%%%%%%%%%%%%%%%%%%%%%%%%%%%%%%%%%%%%%%%%%%%%%

Parallel scalability depends on:

\[
\boxed{
\text{serial fraction},
}
\]

\[
\boxed{
\text{communication volume},
}
\]

\[
\boxed{
\text{synchronization frequency},
}
\]

and

\[
\boxed{
\text{load balance}.
}
\]

Reducing these costs can be as important as increasing processor count.

%%%%%%%%%%%%%%%%%%%%%%%%%%%%%%%%%%%%%%%%%%%%%%%%%%%%%%%%%%%%
\subsection{Performance Portability}
\label{subsec:performance-portability}
%%%%%%%%%%%%%%%%%%%%%%%%%%%%%%%%%%%%%%%%%%%%%%%%%%%%%%%%%%%%

Performance portability means software performs efficiently across
different hardware platforms without requiring a complete redesign.

This is challenging because CPUs, GPUs, and accelerators have different
memory and execution characteristics.

Portable abstractions attempt to separate mathematical algorithms from
hardware-specific details.

%%%%%%%%%%%%%%%%%%%%%%%%%%%%%%%%%%%%%%%%%%%%%%%%%%%%%%%%%%%%
\subsection{Sparse Versus Dense Computation}
\label{subsec:sparse-versus-dense-computation}
%%%%%%%%%%%%%%%%%%%%%%%%%%%%%%%%%%%%%%%%%%%%%%%%%%%%%%%%%%%%

Dense algorithms assume most entries are nonzero.

Sparse algorithms store and process only significant nonzero entries.

If

\[
\operatorname{nnz}(A)
\ll
mn,
\]

then sparse storage can dramatically reduce memory.

However, sparse computations often have lower arithmetic intensity and
more irregular memory access.

%%%%%%%%%%%%%%%%%%%%%%%%%%%%%%%%%%%%%%%%%%%%%%%%%%%%%%%%%%%%
\subsection{Matrix-Free Computation}
\label{subsec:matrix-free-computation}
%%%%%%%%%%%%%%%%%%%%%%%%%%%%%%%%%%%%%%%%%%%%%%%%%%%%%%%%%%%%

A matrix-free method avoids explicitly storing a large matrix.

Instead, it provides an operation

\[
\boxed{
v
\mapsto
Av.
}
\]

Krylov methods often require only matrix-vector products.

For PDEs and optimization, this can eliminate enormous matrix storage
costs.

%%%%%%%%%%%%%%%%%%%%%%%%%%%%%%%%%%%%%%%%%%%%%%%%%%%%%%%%%%%%
\subsection{Matrix-Free Jacobians}
\label{subsec:matrix-free-jacobians}
%%%%%%%%%%%%%%%%%%%%%%%%%%%%%%%%%%%%%%%%%%%%%%%%%%%%%%%%%%%%

A Jacobian-vector product can be approximated by

\[
\boxed{
J(u)v
\approx
\frac{
F(u+\varepsilon v)-F(u)
}{
\varepsilon
}.
}
\]

This permits Newton--Krylov methods without explicitly assembling \(J\).

Automatic differentiation can also compute Jacobian-vector products more
accurately.

%%%%%%%%%%%%%%%%%%%%%%%%%%%%%%%%%%%%%%%%%%%%%%%%%%%%%%%%%%%%
\subsection{Fast Fourier Transform}
\label{subsec:fast-fourier-transform}
%%%%%%%%%%%%%%%%%%%%%%%%%%%%%%%%%%%%%%%%%%%%%%%%%%%%%%%%%%%%

A direct discrete Fourier transform requires approximately

\[
\boxed{
O(N^2)
}
\]

operations.

The Fast Fourier Transform, abbreviated FFT, reduces this to approximately

\[
\boxed{
O(N\log N).
}
\]

This dramatic complexity reduction makes large Fourier-based computations
practical.

%%%%%%%%%%%%%%%%%%%%%%%%%%%%%%%%%%%%%%%%%%%%%%%%%%%%%%%%%%%%
\subsection{Why the FFT Matters}
\label{subsec:why-fft-matters}
%%%%%%%%%%%%%%%%%%%%%%%%%%%%%%%%%%%%%%%%%%%%%%%%%%%%%%%%%%%%

The FFT appears in:

\[
\boxed{
\text{spectral PDE solvers},
}
\]

\[
\boxed{
\text{signal processing},
}
\]

\[
\boxed{
\text{convolution},
}
\]

\[
\boxed{
\text{correlation},
}
\]

and

\[
\boxed{
\text{Poisson solvers}.
}
\]

It is a classic example of algorithmic improvement outweighing raw
hardware speed.

%%%%%%%%%%%%%%%%%%%%%%%%%%%%%%%%%%%%%%%%%%%%%%%%%%%%%%%%%%%%
\subsection{Convolution by FFT}
\label{subsec:fft-convolution}
%%%%%%%%%%%%%%%%%%%%%%%%%%%%%%%%%%%%%%%%%%%%%%%%%%%%%%%%%%%%

The convolution theorem states that convolution in physical space becomes
multiplication in Fourier space.

Conceptually,

\[
\boxed{
f*g
=
\mathcal{F}^{-1}
\left[
\mathcal{F}(f)
\mathcal{F}(g)
\right].
}
\]

Thus large convolutions can be accelerated using FFT algorithms.

%%%%%%%%%%%%%%%%%%%%%%%%%%%%%%%%%%%%%%%%%%%%%%%%%%%%%%%%%%%%
\subsection{Randomized Algorithms}
\label{subsec:scientific-randomized-algorithms}
%%%%%%%%%%%%%%%%%%%%%%%%%%%%%%%%%%%%%%%%%%%%%%%%%%%%%%%%%%%%

Some scientific algorithms use randomness to reduce computational cost.

Examples include:

\[
\boxed{
\text{Monte Carlo simulation},
}
\]

\[
\boxed{
\text{randomized matrix approximation},
}
\]

and

\[
\boxed{
\text{randomized optimization methods}.
}
\]

Randomness becomes part of the numerical algorithm rather than merely a
source of uncertainty.

%%%%%%%%%%%%%%%%%%%%%%%%%%%%%%%%%%%%%%%%%%%%%%%%%%%%%%%%%%%%
\subsection{Randomized Linear Algebra Preview}
\label{subsec:randomized-linear-algebra}
%%%%%%%%%%%%%%%%%%%%%%%%%%%%%%%%%%%%%%%%%%%%%%%%%%%%%%%%%%%%

Suppose \(A\) is a very large matrix.

Instead of computing a complete factorization, one may multiply \(A\) by
a random test matrix

\[
\boxed{
\Omega
}
\]

to form

\[
\boxed{
Y=A\Omega.
}
\]

The range of \(Y\) can approximate the dominant range of \(A\).

This can accelerate low-rank approximation.

%%%%%%%%%%%%%%%%%%%%%%%%%%%%%%%%%%%%%%%%%%%%%%%%%%%%%%%%%%%%
\subsection{Parameter Estimation Pipeline}
\label{subsec:parameter-estimation-pipeline}
%%%%%%%%%%%%%%%%%%%%%%%%%%%%%%%%%%%%%%%%%%%%%%%%%%%%%%%%%%%%

A parameter-estimation problem may contain several nested computational
layers:

\[
\boxed{
\theta
\rightarrow
\text{model simulation}
\rightarrow
\text{predicted data}
\rightarrow
\text{objective}
\rightarrow
\text{optimizer}.
}
\]

Each optimization step may require another simulation.

This can make parameter estimation computationally expensive.

%%%%%%%%%%%%%%%%%%%%%%%%%%%%%%%%%%%%%%%%%%%%%%%%%%%%%%%%%%%%
\subsection{Surrogate Models}
\label{subsec:scientific-surrogate-models}
%%%%%%%%%%%%%%%%%%%%%%%%%%%%%%%%%%%%%%%%%%%%%%%%%%%%%%%%%%%%

If a simulation is extremely expensive, a surrogate model may approximate

\[
\boxed{
\theta
\mapsto
Q(\theta).
}
\]

The surrogate is cheaper to evaluate and can support:

\[
\boxed{
\text{optimization},
}
\]

\[
\boxed{
\text{uncertainty quantification},
}
\]

or

\[
\boxed{
\text{parameter exploration}.
}
\]

The surrogate introduces an additional approximation error that must be
validated.

%%%%%%%%%%%%%%%%%%%%%%%%%%%%%%%%%%%%%%%%%%%%%%%%%%%%%%%%%%%%
\subsection{Reduced-Order Models}
\label{subsec:reduced-order-models}
%%%%%%%%%%%%%%%%%%%%%%%%%%%%%%%%%%%%%%%%%%%%%%%%%%%%%%%%%%%%

A large model with state

\[
u\in\mathbb{R}^N
\]

may be approximated using a lower-dimensional representation

\[
\boxed{
u
\approx
V_r a,
}
\]

where

\[
V_r
\in
\mathbb{R}^{N\times r},
\qquad
r\ll N.
\]

Reduced-order models can accelerate repeated simulations.

%%%%%%%%%%%%%%%%%%%%%%%%%%%%%%%%%%%%%%%%%%%%%%%%%%%%%%%%%%%%
\subsection{Proper Orthogonal Decomposition Preview}
\label{subsec:pod-preview}
%%%%%%%%%%%%%%%%%%%%%%%%%%%%%%%%%%%%%%%%%%%%%%%%%%%%%%%%%%%%

Suppose simulation snapshots form matrix

\[
\boxed{
X
=
\begin{pmatrix}
|&|&&|\\
u_1&u_2&\cdots&u_m\\
|&|&&|
\end{pmatrix}.
}
\]

The singular value decomposition can identify dominant spatial modes.

Keeping the leading modes produces a low-dimensional basis.

This is closely connected with low-rank approximation from
Section~\ref{sec:numerical-linear-algebra}.

%%%%%%%%%%%%%%%%%%%%%%%%%%%%%%%%%%%%%%%%%%%%%%%%%%%%%%%%%%%%
\subsection{Computational Complexity of Simulation}
\label{subsec:simulation-complexity}
%%%%%%%%%%%%%%%%%%%%%%%%%%%%%%%%%%%%%%%%%%%%%%%%%%%%%%%%%%%%

Suppose a time-dependent simulation has:

\[
N
=
\text{spatial unknowns},
\]

and

\[
M
=
\text{time steps}.
\]

If each time step costs

\[
O(N),
\]

then total work is approximately

\[
\boxed{
O(MN).
}
\]

If each step requires a dense solve costing

\[
O(N^3),
\]

the total becomes

\[
\boxed{
O(MN^3).
}
\]

Thus solver choice can completely determine tractability.

%%%%%%%%%%%%%%%%%%%%%%%%%%%%%%%%%%%%%%%%%%%%%%%%%%%%%%%%%%%%
\subsection{Memory Complexity}
\label{subsec:memory-complexity}
%%%%%%%%%%%%%%%%%%%%%%%%%%%%%%%%%%%%%%%%%%%%%%%%%%%%%%%%%%%%

Runtime complexity is not the only resource constraint.

A dense \(N\times N\) matrix requires

\[
\boxed{
O(N^2)
}
\]

storage.

A sparse stencil matrix may require only

\[
\boxed{
O(N)
}
\]

storage.

For very large problems, memory requirements can determine whether a
method is feasible at all.

%%%%%%%%%%%%%%%%%%%%%%%%%%%%%%%%%%%%%%%%%%%%%%%%%%%%%%%%%%%%
\subsection{Time-Memory Tradeoffs}
\label{subsec:time-memory-tradeoffs}
%%%%%%%%%%%%%%%%%%%%%%%%%%%%%%%%%%%%%%%%%%%%%%%%%%%%%%%%%%%%

Some algorithms reduce runtime by storing intermediate quantities.

Others recompute values to reduce memory.

Thus scientific computing often balances:

\[
\boxed{
\text{computation}
}
\]

against

\[
\boxed{
\text{storage}.
}
\]

Checkpointing in adjoint and optimization computations is an important
example.

%%%%%%%%%%%%%%%%%%%%%%%%%%%%%%%%%%%%%%%%%%%%%%%%%%%%%%%%%%%%
\subsection{Parallel Monte Carlo}
\label{subsec:parallel-monte-carlo}
%%%%%%%%%%%%%%%%%%%%%%%%%%%%%%%%%%%%%%%%%%%%%%%%%%%%%%%%%%%%

Suppose

\[
\widehat{\mu}
=
\frac1N
\sum_{i=1}^{N}
g(X_i).
\]

Divide the \(N\) samples across \(p\) workers.

Each computes a partial sum

\[
\boxed{
S_j
=
\sum_{i\in I_j}
g(X_i).
}
\]

Then perform a global reduction:

\[
\boxed{
\widehat{\mu}
=
\frac1N
\sum_{j=1}^{p}
S_j.
}
\]

Only the final reduction requires communication.

%%%%%%%%%%%%%%%%%%%%%%%%%%%%%%%%%%%%%%%%%%%%%%%%%%%%%%%%%%%%
\subsection{Parallel PDE Computation}
\label{subsec:parallel-pde-computation}
%%%%%%%%%%%%%%%%%%%%%%%%%%%%%%%%%%%%%%%%%%%%%%%%%%%%%%%%%%%%

For a stencil-based PDE method, divide the grid into subdomains.

Each processor updates interior points independently.

Only neighboring boundary values must be exchanged.

Thus communication is localized to subdomain interfaces.

The ratio between interior work and boundary communication strongly affects
scaling.

%%%%%%%%%%%%%%%%%%%%%%%%%%%%%%%%%%%%%%%%%%%%%%%%%%%%%%%%%%%%
\subsection{Surface-to-Volume Effect}
\label{subsec:surface-volume-effect}
%%%%%%%%%%%%%%%%%%%%%%%%%%%%%%%%%%%%%%%%%%%%%%%%%%%%%%%%%%%%

In domain decomposition, computation scales with subdomain volume while
communication scales roughly with boundary surface area.

As each processor receives a smaller subdomain, the surface-to-volume ratio
increases.

Thus communication becomes increasingly important in strong scaling.

%%%%%%%%%%%%%%%%%%%%%%%%%%%%%%%%%%%%%%%%%%%%%%%%%%%%%%%%%%%%
\subsection{Parallel Linear Algebra}
\label{subsec:parallel-linear-algebra}
%%%%%%%%%%%%%%%%%%%%%%%%%%%%%%%%%%%%%%%%%%%%%%%%%%%%%%%%%%%%

Dense matrix operations can be distributed across processors by dividing
matrices into blocks.

Sparse iterative methods distribute matrix rows or graph partitions.

The key operations often include:

\[
\boxed{
\text{matrix-vector products},
}
\]

\[
\boxed{
\text{dot products},
}
\]

and

\[
\boxed{
\text{preconditioner applications}.
}
\]

Global dot products can become synchronization bottlenecks.

%%%%%%%%%%%%%%%%%%%%%%%%%%%%%%%%%%%%%%%%%%%%%%%%%%%%%%%%%%%%
\subsection{Asynchronous Computation}
\label{subsec:asynchronous-computation}
%%%%%%%%%%%%%%%%%%%%%%%%%%%%%%%%%%%%%%%%%%%%%%%%%%%%%%%%%%%%

Traditional parallel algorithms often require all processors to synchronize
at specific points.

Asynchronous methods allow workers to proceed using partially outdated
information.

This can reduce waiting but complicates convergence analysis.

Some iterative optimization and linear-solver methods admit asynchronous
variants.

%%%%%%%%%%%%%%%%%%%%%%%%%%%%%%%%%%%%%%%%%%%%%%%%%%%%%%%%%%%%
\subsection{Fault Tolerance}
\label{subsec:scientific-fault-tolerance}
%%%%%%%%%%%%%%%%%%%%%%%%%%%%%%%%%%%%%%%%%%%%%%%%%%%%%%%%%%%%

As computational systems grow, hardware or software failures become more
likely during long simulations.

Fault-tolerance strategies include:

\[
\boxed{
\text{checkpoint/restart},
}
\]

\[
\boxed{
\text{redundancy},
}
\]

and

\[
\boxed{
\text{algorithmic recovery}.
}
\]

Scientific workflows should anticipate failure rather than assume every
long run will finish uninterrupted.

%%%%%%%%%%%%%%%%%%%%%%%%%%%%%%%%%%%%%%%%%%%%%%%%%%%%%%%%%%%%
\subsection{Scientific Computing and Cloud Systems}
\label{subsec:scientific-cloud-computing}
%%%%%%%%%%%%%%%%%%%%%%%%%%%%%%%%%%%%%%%%%%%%%%%%%%%%%%%%%%%%

Scientific workloads may run on local workstations, clusters,
supercomputers, or cloud-based systems.

The mathematical algorithms remain similar, but resource management and
cost models differ.

Cloud computation can provide flexible access to:

\[
\boxed{
\text{CPUs},
}
\]

\[
\boxed{
\text{GPUs},
}
\]

and

\[
\boxed{
\text{distributed storage}.
}
\]

%%%%%%%%%%%%%%%%%%%%%%%%%%%%%%%%%%%%%%%%%%%%%%%%%%%%%%%%%%%%
\subsection{Reproducibility Across Hardware}
\label{subsec:reproducibility-across-hardware}
%%%%%%%%%%%%%%%%%%%%%%%%%%%%%%%%%%%%%%%%%%%%%%%%%%%%%%%%%%%%

Changing hardware can alter:

\[
\boxed{
\text{floating-point reduction order},
}
\]

\[
\boxed{
\text{compiler optimizations},
}
\]

and

\[
\boxed{
\text{mathematical-library implementations}.
}
\]

Therefore numerical agreement should often be assessed using scientifically
meaningful tolerances rather than demanding exact identical digits.

%%%%%%%%%%%%%%%%%%%%%%%%%%%%%%%%%%%%%%%%%%%%%%%%%%%%%%%%%%%%
\subsection{Scientific Computing Workflow for Large Problems}
\label{subsec:large-scientific-computing-workflow}
%%%%%%%%%%%%%%%%%%%%%%%%%%%%%%%%%%%%%%%%%%%%%%%%%%%%%%%%%%%%

For a large computational project:

\begin{enumerate}
    \item begin with a small verified model;
    \item compare against analytic or benchmark results;
    \item perform a discretization study;
    \item establish numerical tolerances;
    \item profile the implementation;
    \item identify dominant computational kernels;
    \item exploit efficient libraries;
    \item exploit sparsity;
    \item reduce unnecessary data movement;
    \item parallelize only meaningful bottlenecks;
    \item measure strong and weak scaling;
    \item implement checkpointing;
    \item document all runtime parameters;
    \item automate parameter sweeps and postprocessing;
    \item verify that high-performance changes preserve numerical results;
    \item perform final scientific validation.
\end{enumerate}

%%%%%%%%%%%%%%%%%%%%%%%%%%%%%%%%%%%%%%%%%%%%%%%%%%%%%%%%%%%%
\subsection{Choosing Computational Strategies}
\label{subsec:choosing-computational-strategies}
%%%%%%%%%%%%%%%%%%%%%%%%%%%%%%%%%%%%%%%%%%%%%%%%%%%%%%%%%%%%

For small dense linear algebra, use optimized dense factorizations.

For large sparse systems, use:

\[
\boxed{
\text{sparse direct methods}
}
\]

or

\[
\boxed{
\text{preconditioned iterative methods}.
}
\]

For independent simulation ensembles, use:

\[
\boxed{
\text{task parallelism}.
}
\]

For stencil-based PDEs, use:

\[
\boxed{
\text{domain decomposition}.
}
\]

For large data-parallel kernels, consider:

\[
\boxed{
\text{GPU acceleration}.
}
\]

For memory-limited systems, consider:

\[
\boxed{
\text{sparse or matrix-free algorithms}.
}
\]

%%%%%%%%%%%%%%%%%%%%%%%%%%%%%%%%%%%%%%%%%%%%%%%%%%%%%%%%%%%%
\subsection{Common Errors}
\label{subsec:scientific-computing-common-errors}
%%%%%%%%%%%%%%%%%%%%%%%%%%%%%%%%%%%%%%%%%%%%%%%%%%%%%%%%%%%%

\subsubsection{Optimizing Code Before Verifying Mathematics}

A fast implementation of an incorrect method is still incorrect.

Correctness should be established before performance optimization.

%%%%%%%%%%%%%%%%%%%%%%%%%%%%%%%%%%%%%%%%%%%%%%%%%%%%%%%%%%%%
\subsubsection{Assuming Big-O Complexity Predicts Exact Runtime}

Two algorithms with the same asymptotic complexity can have very different
constants, memory behavior, and hardware utilization.

Benchmark representative problems.

%%%%%%%%%%%%%%%%%%%%%%%%%%%%%%%%%%%%%%%%%%%%%%%%%%%%%%%%%%%%
\subsubsection{Ignoring Memory Access}

Modern computational performance may be limited more by data movement than
by arithmetic.

Reducing operation count alone may not improve runtime.

%%%%%%%%%%%%%%%%%%%%%%%%%%%%%%%%%%%%%%%%%%%%%%%%%%%%%%%%%%%%
\subsubsection{Reimplementing Standard Numerical Kernels Unnecessarily}

Fundamental matrix and transform operations are often already available in
highly optimized and well-tested libraries.

Use established kernels unless a specialized implementation is required.

%%%%%%%%%%%%%%%%%%%%%%%%%%%%%%%%%%%%%%%%%%%%%%%%%%%%%%%%%%%%
\subsubsection{Parallelizing a Tiny Problem}

Parallel overhead can exceed the useful computation.

Parallelism becomes beneficial only when sufficient work exists to justify
communication and scheduling costs.

%%%%%%%%%%%%%%%%%%%%%%%%%%%%%%%%%%%%%%%%%%%%%%%%%%%%%%%%%%%%
\subsubsection{Assuming Twice as Many Processors Gives Twice the Speed}

Serial work, communication, and synchronization prevent ideal scaling.

Measure speedup directly.

%%%%%%%%%%%%%%%%%%%%%%%%%%%%%%%%%%%%%%%%%%%%%%%%%%%%%%%%%%%%
\subsubsection{Ignoring Load Imbalance}

Parallel runtime is often determined by the slowest worker.

Equal processor counts do not imply equal workloads.

%%%%%%%%%%%%%%%%%%%%%%%%%%%%%%%%%%%%%%%%%%%%%%%%%%%%%%%%%%%%
\subsubsection{Using Many Small Network Messages}

Latency costs can dominate.

Combining communication into fewer larger messages may improve
performance.

%%%%%%%%%%%%%%%%%%%%%%%%%%%%%%%%%%%%%%%%%%%%%%%%%%%%%%%%%%%%
\subsubsection{Transferring Accelerator Data Repeatedly}

GPU computation can become slower than CPU computation if data movement
dominates kernel execution.

Keep frequently used data close to the accelerator when possible.

%%%%%%%%%%%%%%%%%%%%%%%%%%%%%%%%%%%%%%%%%%%%%%%%%%%%%%%%%%%%
\subsubsection{Assuming Lower Precision Is Always Safe}

Lower precision can increase rounding error or cause convergence failure
for ill-conditioned problems.

Precision should be selected using numerical analysis.

%%%%%%%%%%%%%%%%%%%%%%%%%%%%%%%%%%%%%%%%%%%%%%%%%%%%%%%%%%%%
\subsubsection{Assuming Double Precision Is Always Necessary}

Some well-conditioned calculations can be performed accurately in lower
precision, potentially reducing cost.

Mixed precision may provide a better tradeoff.

%%%%%%%%%%%%%%%%%%%%%%%%%%%%%%%%%%%%%%%%%%%%%%%%%%%%%%%%%%%%
\subsubsection{Expecting Parallel Floating-Point Reductions to Be Bitwise Identical}

Different summation orders can change low-order digits.

This is normal floating-point behavior and should be distinguished from
meaningful numerical disagreement.

%%%%%%%%%%%%%%%%%%%%%%%%%%%%%%%%%%%%%%%%%%%%%%%%%%%%%%%%%%%%
\subsubsection{Using Random Simulation Without Recording Seeds}

Exact reproduction of stochastic experiments becomes difficult.

Record both generator configuration and seed information.

%%%%%%%%%%%%%%%%%%%%%%%%%%%%%%%%%%%%%%%%%%%%%%%%%%%%%%%%%%%%
\subsubsection{Saving Too Little Information to Reproduce Results}

A final graph alone is not a reproducible computational experiment.

Preserve model parameters, code version, numerical settings, and input
data.

%%%%%%%%%%%%%%%%%%%%%%%%%%%%%%%%%%%%%%%%%%%%%%%%%%%%%%%%%%%%
\subsubsection{Saving Every Possible Output}

Massive unnecessary output can dominate runtime and storage.

Save quantities needed for scientific interpretation and reproducibility.

%%%%%%%%%%%%%%%%%%%%%%%%%%%%%%%%%%%%%%%%%%%%%%%%%%%%%%%%%%%%
\subsubsection{Treating a Plot as Verification}

A smooth-looking solution can be quantitatively wrong.

Use residuals, convergence studies, conservation checks, and benchmark
solutions.

%%%%%%%%%%%%%%%%%%%%%%%%%%%%%%%%%%%%%%%%%%%%%%%%%%%%%%%%%%%%
\subsubsection{Ignoring Units}

A unit mismatch can produce numerically plausible but physically
meaningless results.

Perform dimensional checks.

%%%%%%%%%%%%%%%%%%%%%%%%%%%%%%%%%%%%%%%%%%%%%%%%%%%%%%%%%%%%
\subsubsection{Failing to Test Small Problems First}

Large simulations are difficult to debug.

Start with small cases where solutions can be inspected or compared with
known results.

%%%%%%%%%%%%%%%%%%%%%%%%%%%%%%%%%%%%%%%%%%%%%%%%%%%%%%%%%%%%
\subsubsection{Changing Algorithm and Hardware Simultaneously}

If performance or results change, it becomes difficult to identify the
cause.

Controlled computational experiments change one major factor at a time.

%%%%%%%%%%%%%%%%%%%%%%%%%%%%%%%%%%%%%%%%%%%%%%%%%%%%%%%%%%%%
\subsubsection{Assuming More Hardware Fixes a Poor Algorithm}

An algorithm with poor asymptotic scaling may remain inefficient even on a
large computer.

Algorithmic improvements often provide much larger gains.

%%%%%%%%%%%%%%%%%%%%%%%%%%%%%%%%%%%%%%%%%%%%%%%%%%%%%%%%%%%%
\subsection{Applications}
\label{subsec:scientific-computing-applications}
%%%%%%%%%%%%%%%%%%%%%%%%%%%%%%%%%%%%%%%%%%%%%%%%%%%%%%%%%%%%

\begin{applicationbox}

\textbf{Physics.}
Scientific computing supports particle simulations, quantum models,
electromagnetics, statistical mechanics, and astrophysics.

\medskip

\textbf{Engineering.}
Large numerical simulations are used in structures, fluids, heat transfer,
circuits, control, and mechanical design.

\medskip

\textbf{Biology.}
Computational models study epidemics, ecological systems, molecular
systems, and population dynamics.

\medskip

\textbf{Chemistry.}
Reaction kinetics, molecular simulation, quantum chemistry, and transport
models rely on high-performance numerical computation.

\medskip

\textbf{Climate and Environment.}
Atmospheric, oceanic, and environmental models involve extremely large
coupled PDE systems and ensemble simulations.

\medskip

\textbf{Statistics and Data Science.}
Large-scale linear algebra, optimization, Bayesian computation, and
simulation require efficient numerical software.

\medskip

\textbf{Finance.}
Monte Carlo simulation, PDE methods, optimization, and stochastic models
are used in pricing and risk analysis.

\medskip

\textbf{Operations Research.}
Large optimization, simulation, and network models depend on efficient
computational algorithms.

\end{applicationbox}

%%%%%%%%%%%%%%%%%%%%%%%%%%%%%%%%%%%%%%%%%%%%%%%%%%%%%%%%%%%%
\subsection{Exercises}
\label{subsec:scientific-computing-exercises}
%%%%%%%%%%%%%%%%%%%%%%%%%%%%%%%%%%%%%%%%%%%%%%%%%%%%%%%%%%%%

\begin{exercise}[1]
Define scientific computing.
\end{exercise}

\begin{exercise}[1]
Distinguish a mathematical model from a numerical method.
\end{exercise}

\begin{exercise}[1]
Define algorithmic complexity.
\end{exercise}

\begin{exercise}[1]
Define arithmetic intensity conceptually.
\end{exercise}

\begin{exercise}[1]
Define parallel speedup.
\end{exercise}

\begin{exercise}[1]
Define parallel efficiency.
\end{exercise}

\begin{exercise}[1]
Define strong scaling.
\end{exercise}

\begin{exercise}[1]
Define weak scaling.
\end{exercise}

\begin{exercise}[1]
Define checkpointing.
\end{exercise}

\begin{exercise}[1]
Define data provenance.
\end{exercise}

\begin{exercise}[2]
Compare the operation counts of algorithms requiring

\[
O(n),
\qquad
O(n^2),
\qquad
O(n^3)
\]

work as \(n\) increases.
\end{exercise}

\begin{exercise}[2]
Explain why a sparse matrix-vector product may be memory-bound.
\end{exercise}

\begin{exercise}[2]
Explain temporal and spatial locality.
\end{exercise}

\begin{exercise}[2]
Give an example of a vectorized mathematical operation.
\end{exercise}

\begin{exercise}[2]
Explain why numerical libraries can outperform scalar high-level loops.
\end{exercise}

\begin{exercise}[2]
Compute speedup from given values of \(T_1\) and \(T_p\).
\end{exercise}

\begin{exercise}[2]
Compute parallel efficiency.
\end{exercise}

\begin{exercise}[3]
Apply Amdahl's law for a specified serial fraction and processor count.
\end{exercise}

\begin{exercise}[3]
Compute the maximum theoretical speedup as

\[
p\to\infty.
\]
\end{exercise}

\begin{exercise}[3]
Explain why strong-scaling efficiency eventually decreases.
\end{exercise}

\begin{exercise}[3]
Design a weak-scaling experiment.
\end{exercise}

\begin{exercise}[3]
Give an example of task parallelism.
\end{exercise}

\begin{exercise}[3]
Give an example of data parallelism.
\end{exercise}

\begin{exercise}[3]
Explain the difference between shared and distributed memory.
\end{exercise}

\begin{exercise}[3]
Explain a race condition.
\end{exercise}

\begin{exercise}[3]
Explain why synchronization reduces parallel efficiency.
\end{exercise}

\begin{exercise}[3]
Describe halo exchange for a finite-difference PDE solver.
\end{exercise}

\begin{exercise}[3]
Explain why Monte Carlo is naturally parallel.
\end{exercise}

\begin{exercise}[3]
Explain why global dot products can become bottlenecks in Krylov solvers.
\end{exercise}

\begin{exercise}[3]
Compare single and double precision conceptually.
\end{exercise}

\begin{exercise}[3]
Describe a mixed-precision linear-solver strategy.
\end{exercise}

\begin{exercise}[3]
Explain why parallel floating-point sums can differ in their final digits.
\end{exercise}

\begin{exercise}[3]
Compare sequential and pairwise summation.
\end{exercise}

\begin{exercise}[3]
Explain compensated summation.
\end{exercise}

\begin{exercise}[3]
Distinguish bitwise and numerical reproducibility.
\end{exercise}

\begin{exercise}[3]
List the information required to reproduce a computational experiment.
\end{exercise}

\begin{exercise}[3]
Design a checkpoint strategy for a twelve-hour simulation.
\end{exercise}

\begin{exercise}[3]
Explain the tradeoff in checkpoint frequency.
\end{exercise}

\begin{exercise}[3]
Construct a parameter sweep over two model parameters.
\end{exercise}

\begin{exercise}[3]
Explain why parameter sweeps are often embarrassingly parallel.
\end{exercise}

\begin{exercise}[3]
Design an ensemble simulation.
\end{exercise}

\begin{exercise}[3]
Construct a simple computational dependency graph.
\end{exercise}

\begin{exercise}[3]
Explain the difference between unit testing and regression testing.
\end{exercise}

\begin{exercise}[3]
Design a convergence test for a second-order numerical method.
\end{exercise}

\begin{exercise}[3]
Explain why a small residual does not guarantee a small discretization
error.
\end{exercise}

\begin{exercise}[3]
Construct a conservation diagnostic for a population model.
\end{exercise}

\begin{exercise}[3]
Nondimensionalize a simple ODE using a characteristic time scale.
\end{exercise}

\begin{exercise}[3]
Explain the difference between sensitivity and uncertainty.
\end{exercise}

\begin{exercise}[3]
Design a quantitative diagnostic to accompany a scientific visualization.
\end{exercise}

\begin{exercise}[3]
Explain why excessive output can reduce simulation performance.
\end{exercise}

\begin{exercise}[3]
Compare a dense \(N\times N\) storage requirement with a sparse matrix
containing \(O(N)\) nonzeros.
\end{exercise}

\begin{exercise}[3]
Explain matrix-free computation.
\end{exercise}

\begin{exercise}[3]
Construct a finite-difference Jacobian-vector product.
\end{exercise}

\begin{exercise}[3]
Compare the asymptotic costs of a direct DFT and an FFT.
\end{exercise}

\begin{exercise}[3]
Explain why the FFT is important in spectral PDE solvers.
\end{exercise}

\begin{exercise}[3]
Describe a randomized low-rank matrix approximation.
\end{exercise}

\begin{exercise}[3]
Explain the role of surrogate models in expensive optimization.
\end{exercise}

\begin{exercise}[3]
Describe a reduced-order representation

\[
u\approx V_ra.
\]
\end{exercise}

\begin{exercise}[3]
Estimate the computational cost of a simulation with \(M\) time steps and
\(O(N)\) work per step.
\end{exercise}

\begin{exercise}[3]
Explain the surface-to-volume effect in parallel PDE decomposition.
\end{exercise}

\begin{exercise}[3]
Explain the purpose of profiling.
\end{exercise}

\begin{exercise}[3]
Design a fair benchmark comparing two numerical algorithms.
\end{exercise}

\begin{exercise}[3]
Explain why changing problem size can change which algorithm is fastest.
\end{exercise}

\begin{exercise}[4]
Study the roofline performance model.
\end{exercise}

\begin{exercise}[4]
Derive Amdahl's law.
\end{exercise}

\begin{exercise}[4]
Study Gustafson-type scaled-speedup arguments.
\end{exercise}

\begin{exercise}[4]
Analyze communication complexity of a domain-decomposed PDE method.
\end{exercise}

\begin{exercise}[4]
Study parallel matrix multiplication.
\end{exercise}

\begin{exercise}[4]
Study communication-avoiding linear algebra.
\end{exercise}

\begin{exercise}[4]
Investigate parallel Krylov methods.
\end{exercise}

\begin{exercise}[4]
Study asynchronous iterative methods.
\end{exercise}

\begin{exercise}[4]
Analyze load balancing for adaptive meshes.
\end{exercise}

\begin{exercise}[4]
Study mixed-precision iterative refinement.
\end{exercise}

\begin{exercise}[4]
Analyze floating-point reproducibility of parallel reductions.
\end{exercise}

\begin{exercise}[4]
Study compensated and reproducible summation algorithms.
\end{exercise}

\begin{exercise}[4]
Study arbitrary-precision computational complexity.
\end{exercise}

\begin{exercise}[4]
Study checkpointing models for failure-prone large-scale systems.
\end{exercise}

\begin{exercise}[4]
Study parallel pseudorandom number generation.
\end{exercise}

\begin{exercise}[4]
Investigate randomized numerical linear algebra.
\end{exercise}

\begin{exercise}[4]
Study matrix-free Newton--Krylov methods.
\end{exercise}

\begin{exercise}[4]
Study proper orthogonal decomposition for reduced-order modeling.
\end{exercise}

\begin{exercise}[4]
Investigate in situ visualization and analysis.
\end{exercise}

\begin{exercise}[4]
Study performance portability across CPUs and GPUs.
\end{exercise}

\begin{exercise}[4]
Design a fully reproducible large-scale computational experiment.
\end{exercise}

%%%%%%%%%%%%%%%%%%%%%%%%%%%%%%%%%%%%%%%%%%%%%%%%%%%%%%%%%%%%
\subsection{Conceptual Summary}
\label{subsec:scientific-computing-summary}
%%%%%%%%%%%%%%%%%%%%%%%%%%%%%%%%%%%%%%%%%%%%%%%%%%%%%%%%%%%%

Scientific computing combines

\[
\boxed{
\text{mathematical models}
+
\text{numerical algorithms}
+
\text{software}
+
\text{computer architecture}.
}
\]

Computational cost depends on more than arithmetic.

It also depends on

\[
\boxed{
\text{memory movement}
+
\text{communication}
+
\text{synchronization}.
}
\]

Parallel speedup is

\[
\boxed{
S_p
=
\frac{
T_1
}{
T_p
},
}
\]

and parallel efficiency is

\[
\boxed{
E_p
=
\frac{
S_p
}{
p
}.
}
\]

Amdahl's law gives

\[
\boxed{
S_p
=
\frac{
1
}{
s+\frac{1-s}{p}
}.
}
\]

Large scientific calculations rely on

\[
\boxed{
\text{sparse algorithms},
}
\]

\[
\boxed{
\text{parallelism},
}
\]

\[
\boxed{
\text{accelerators},
}
\]

\[
\boxed{
\text{checkpointing},
}
\]

and

\[
\boxed{
\text{reproducible workflows}.
}
\]

Reliable scientific computing therefore requires simultaneous attention to

\[
\boxed{
\text{correct mathematics}
+
\text{numerical stability}
+
\text{computational efficiency}
+
\text{software verification}
+
\text{scientific reproducibility}.
}
\]

%%%%%%%%%%%%%%%%%%%%%%%%%%%%%%%%%%%%%%%%%%%%%%%%%%%%%%%%%%%%
\subsection{Section Connections}
\label{subsec:scientific-computing-connections}
%%%%%%%%%%%%%%%%%%%%%%%%%%%%%%%%%%%%%%%%%%%%%%%%%%%%%%%%%%%%

\chapterconnection{
Scientific Computing integrates the numerical methods developed throughout
Chapter~11 into complete computational workflows. Numerical Analysis
provides error and stability principles; Numerical Linear Algebra supplies
matrix algorithms; Numerical Optimization supplies iterative decision
methods; Numerical ODEs and PDEs generate large simulation problems; and
Monte Carlo Methods supply stochastic computation. Scientific Computing
adds the implementation principles required to execute these methods
reliably and efficiently, including memory-aware algorithms, parallelism,
high-performance computing, precision selection, checkpointing,
verification, and reproducibility. The next section, Symbolic Computation,
moves from approximate numerical calculation to exact manipulation of
mathematical expressions, algebraic identities, equations, derivatives,
integrals, and symbolic mathematical objects.
}

%%%%%%%%%%%%%%%%%%%%%%%%%%%%%%%%%%%%%%%%%%%%%%%%%%%%%%%%%%%%
% SECTION SOURCE TRACKING
%%%%%%%%%%%%%%%%%%%%%%%%%%%%%%%%%%%%%%%%%%%%%%%%%%%%%%%%%%%%

%------------------------------------------------------------
% 11.7 Scientific Computing
%------------------------------------------------------------
%
% Source basis:
%   - Standard scientific-computing principles.
%   - Standard high-performance and parallel-computing concepts.
%   - Standard numerical-software verification and reproducibility
%     practices.
%
% Used for:
%   - scientific computing definition;
%   - model, algorithm, and implementation separation;
%   - computational error sources;
%   - algorithmic complexity;
%   - operation counts;
%   - memory hierarchy;
%   - locality;
%   - arithmetic intensity;
%   - memory-bound and compute-bound computation;
%   - vectorization;
%   - numerical libraries;
%   - profiling and benchmarking;
%   - parallel computing;
%   - task and data parallelism;
%   - shared-memory systems;
%   - race conditions and synchronization;
%   - distributed-memory systems;
%   - message passing;
%   - domain decomposition;
%   - halo exchange;
%   - speedup and parallel efficiency;
%   - Amdahl's law;
%   - strong and weak scaling;
%   - load balancing;
%   - communication overhead;
%   - latency and bandwidth;
%   - granularity;
%   - embarrassingly parallel computation;
%   - GPU computing;
%   - accelerator data movement;
%   - single, double, mixed, and arbitrary precision;
%   - floating-point reduction order;
%   - pairwise and compensated summation;
%   - deterministic and nondeterministic computation;
%   - bitwise and numerical reproducibility;
%   - computational experiment documentation;
%   - data provenance;
%   - configuration and logging;
%   - checkpointing;
%   - parameter sweeps;
%   - ensemble simulation;
%   - batch processing;
%   - workflow automation;
%   - unit and regression testing;
%   - convergence and residual testing;
%   - conservation checks;
%   - dimensional analysis;
%   - nondimensionalization;
%   - sensitivity studies;
%   - scientific visualization;
%   - reduction operations;
%   - I/O and in situ analysis;
%   - high-performance computing;
%   - scalability;
%   - performance portability;
%   - sparse and matrix-free computation;
%   - Jacobian-vector products;
%   - Fast Fourier Transform;
%   - randomized algorithms;
%   - randomized linear algebra preview;
%   - parameter-estimation pipelines;
%   - surrogate models;
%   - reduced-order models;
%   - proper orthogonal decomposition preview;
%   - simulation and memory complexity;
%   - parallel Monte Carlo;
%   - parallel PDE computation;
%   - surface-to-volume scaling;
%   - parallel linear algebra;
%   - asynchronous computation;
%   - fault tolerance;
%   - applications and exercises.
%
% Notes:
%   - Symbolic Computation is developed next in Section 11.8.
%   - Computer Algebra follows in Section 11.9.
%   - Mathematical Simulation follows in Section 11.10.
%   - High-Performance Mathematical Computing is developed in
%     Section 11.11.
%
%%%%%%%%%%%%%%%%%%%%%%%%%%%%%%%%%%%%%%%%%%%%%%%%%%%%%%%%%%%%